\documentclass[../DoAn.tex]{subfiles}
\begin{document}

Chương này đi từ yêu cầu tới thiết kế. Mục~\ref{section:3.1} đặt các tác nhân, ca
sử dụng và yêu cầu, mỗi yêu cầu mang một mã để các mục sau viện dẫn lại. Các mục
còn lại trình bày kiến trúc tổng thể, phần cứng, pipeline AI, firmware, bản tin
MQTT, cơ sở dữ liệu, API và giao diện. Chương này nói hệ thống gồm những gì và vì
sao, còn cách làm nằm ở Chương~\ref{chapter:Implementation}.

\section{Phân tích yêu cầu}
\label{section:3.1}

\subsection{Tác nhân và ca sử dụng}
\label{section:3.1.1}

Hệ thống có năm nhóm người dùng và một tác nhân là máy chấm công, liệt kê ở
Bảng~\ref{tab:tac_nhan}. Vai quản lý không ai giao tay mà tự tính từ cây báo cáo: ai
có người dưới quyền thì mang vai đó. Tài khoản chưa ai giao việc mang vai thứ sáu,
\texttt{VIEWER}, và chỉ thấy chính mình.

\begin{table}[!htb]
\centering
\small
\caption{Các tác nhân của hệ thống.}
\label{tab:tac_nhan}
\begin{tabular}{|L{3.2cm}|L{2.4cm}|L{7.9cm}|}
\hline
\textbf{Tác nhân} & \textbf{Vai} & \textbf{Việc chính} \\
\hline
Nhân viên & \texttt{EMPLOYEE} & Chấm công tại máy, xem công, phép, phiếu lương của mình, gửi đơn \\
\hline
Quản lý & \texttt{MANAGER} & Duyệt đơn và xem người trong cây dưới quyền, không xem lương và giấy tờ riêng của họ \\
\hline
Nhân sự & \texttt{HR} & Hồ sơ, hợp đồng, ca làm, bảng công, phép, mở lượt đăng ký khuôn mặt \\
\hline
Kế toán lương & \texttt{PAYROLL} & Kỳ lương, chốt kỳ, phiếu lương \\
\hline
Quản trị & \texttt{ADMIN} & Máy chấm công, bản phát hành, tài khoản, nhật ký \\
\hline
Máy chấm công & Không có & Gửi lượt chấm và trạng thái, nhận lệnh, danh sách người và bản cập nhật \\
\hline
\end{tabular}
\end{table}

Hình~\ref{fig:use_case} là sơ đồ ca sử dụng tổng quát. Bên trái là việc của nhân viên,
quản lý và nhân sự, bên phải là việc của kế toán lương và quản trị. Ca đăng ký khuôn
mặt có hai tác nhân vì nhân sự mở lượt còn nhân viên đứng trước máy. Ca duyệt đơn có hai
tác nhân vì đơn đi tới quản lý hoặc nhân sự.

\newcommand{\actor}[3]{%
  \begin{scope}[shift={(#1)}]
    \draw (0,0.42) circle (0.13);
    \draw (0,0.29) -- (0,-0.08);
    \draw (-0.2,0.16) -- (0.2,0.16);
    \draw (0,-0.08) -- (-0.15,-0.36);
    \draw (0,-0.08) -- (0.15,-0.36);
    \node[below, align=center] at (0,-0.36) {#3};
    \coordinate (#2) at (0,0.1);
  \end{scope}}

\tikzset{
  sqsync/.style={-{Triangle[length=1.8mm, width=1.4mm]}},
  sqasync/.style={-{Straight Barb[length=1.8mm, width=2mm]}},
  sqreply/.style={densely dashed, -{Straight Barb[length=1.8mm, width=2mm]}},
  sqlab/.style={font=\scriptsize, above, align=center, fill=white, inner sep=1.5pt},
  sqnote/.style={font=\scriptsize, fill=white, inner sep=1pt},
  incl/.style={densely dashed, -{Straight Barb[length=1.8mm, width=2mm]}},
  erfk/.style={preaction={draw, arrows=-}, {Straight Barb[reversed, length=1.6mm, width=2.6mm] Circle[open, fill=white, length=1.2mm, sep=0.6mm]}-{Bar[width=2.4mm, sep=0.8mm] Bar[width=2.4mm]}},
  erfkopt/.style={preaction={draw, arrows=-}, {Straight Barb[reversed, length=1.6mm, width=2.6mm] Circle[open, fill=white, length=1.2mm, sep=0.6mm]}-{Circle[open, fill=white, length=1.2mm, sep=0.6mm] Bar[width=2.4mm]}},
  eroo/.style={{Bar[width=2.4mm, sep=0.6mm] Circle[open, fill=white, length=1.2mm]}-{Circle[open, fill=white, length=1.2mm, sep=0.6mm] Bar[width=2.4mm]}}}
\newcommand{\sqobj}[4]{%
  \node[draw, fill=white, text width=#2, align=center, minimum height=7mm, inner sep=1pt] at (#1,0) {#3};
  \draw[densely dashed] (#1,-0.35) -- (#1,#4);}
\newcommand{\sqactor}[3]{%
  \actor{#1,0.25}{sqa}{#2}
  \draw[densely dashed] (#1,-0.55) -- (#1,#3);}
\newcommand{\sqexec}[3]{\draw[fill=white] ({#1-0.08},#2) rectangle ({#1+0.08},#3);}
\newcommand{\sqselfr}[4][1]{%
  \ifnum#1=1 \sqexec{#2}{#3+0.1}{#3-0.42}\fi
  \draw[sqsync] ({#2+0.08},#3) -- ({#2+0.5},#3) -- ({#2+0.5},{#3-0.32}) -- ({#2+0.08},{#3-0.32});
  \node[sqnote, anchor=west, align=left] at ({#2+0.55},{#3-0.16}) {#4};}
\newcommand{\sqselfl}[4][1]{%
  \ifnum#1=1 \sqexec{#2}{#3+0.1}{#3-0.42}\fi
  \draw[sqsync] ({#2-0.08},#3) -- ({#2-0.5},#3) -- ({#2-0.5},{#3-0.32}) -- ({#2-0.08},{#3-0.32});
  \node[sqnote, anchor=east, align=right] at ({#2-0.55},{#3-0.16}) {#4};}
\newcommand{\sqfrag}[5]{%
  \draw (#1,#2) rectangle (#3,#4);
  \node[font=\scriptsize, anchor=north west, inner sep=2pt] (sqf) at (#1,#2) {#5};
  \draw (sqf.south west) -- ([xshift=-1.2mm]sqf.south east) -- ([yshift=1.2mm]sqf.south east) -- (sqf.north east);}
\newcommand{\bcehead}[3]{%
  \draw (#1,0.15) circle (0.22);
  \node[font=\scriptsize, below, align=center, inner sep=1pt] at (#1,-0.1) {#2};
  \draw[densely dashed] (#1,-0.85) -- (#1,#3);}
\newcommand{\bceboundary}[3]{%
  \bcehead{#1}{#2}{#3}
  \draw ({#1-0.42},-0.05) -- ({#1-0.42},0.35);
  \draw ({#1-0.42},0.15) -- ({#1-0.22},0.15);}
\newcommand{\bcecontrol}[3]{%
  \bcehead{#1}{#2}{#3}
  \draw ({#1+0.08},0.45) -- ({#1-0.03},0.37) -- ({#1+0.08},0.29);}
\newcommand{\bceentity}[3]{%
  \bcehead{#1}{#2}{#3}
  \draw ({#1-0.22},-0.07) -- ({#1+0.22},-0.07);}

\begin{figure}[!htb]
\centering
\begin{tikzpicture}[
  font=\footnotesize,
  uc/.style={draw, ellipse, align=center, minimum width=33mm, minimum height=9mm, inner sep=0pt}]
\draw (-3.5,-3.8) rectangle (3.95,4.0);
\node[anchor=north west] at (-3.45,3.97) {Hệ thống chấm công};
\node[uc] (a1) at (-1.65,2.85) {Chấm công tại máy};
\node[uc] (a3) at (-1.65,1.65) {Gửi đơn nghỉ,\\đơn làm thêm};
\node[uc] (a2) at (-1.65,0.45) {Xem công, phép,\\phiếu lương};
\node[uc] (a5) at (-1.65,-0.75) {Đăng ký khuôn mặt\\tại máy};
\node[uc] (a4) at (-1.65,-1.95) {Duyệt đơn};
\node[uc] (b1) at (-1.65,-3.15) {Quản lý hồ sơ,\\ca, bảng công};
\node[uc] (b2) at (2.15,2.1) {Chạy và chốt\\kỳ lương};
\node[uc] (b3) at (2.15,0.5) {Duyệt máy mới};
\node[uc] (b4) at (2.15,-1.0) {Theo dõi, ra lệnh\\cho máy};
\node[uc] (b5) at (2.15,-2.5) {Phát hành bản\\cập nhật OTA};
\actor{-6.1,2.0}{e}{Nhân viên}
\actor{-6.1,-0.4}{m}{Quản lý}
\actor{-6.1,-2.8}{h}{Nhân sự}
\actor{6.2,2.0}{p}{Kế toán lương}
\actor{6.2,-1.1}{q}{Quản trị}
\draw (e) -- (a1.west);
\draw (e) -- (a2.west);
\draw (e) -- (a3.west);
\draw (e) -- (a5.west);
\draw (m) -- (a2.west);
\draw (m) -- (a4.west);
\draw (h) -- (a4.west);
\draw (h) -- (a5.west);
\draw (h) -- (b1.west);
\draw (p) -- (b2.east);
\draw (q) -- (b3.east);
\draw (q) -- (b4.east);
\draw (q) -- (b5.east);
\end{tikzpicture}
\caption{Sơ đồ ca sử dụng tổng quát.}
\label{fig:use_case}
\end{figure}

Hình~\ref{fig:uc_may}, Hình~\ref{fig:uc_ban_giay} và Hình~\ref{fig:uc_nhan_vien} phân rã sơ đồ tổng quát theo ba nhóm việc. Ở phần
máy chấm công, chấm công và đăng ký khuôn mặt cùng gọi ca xác thực khuôn mặt, vì đăng ký đi qua đúng đường phát hiện và chống giả mạo
của chấm công. Máy chấm công là tác nhân của các ca do chính nó mở: xin duyệt ở lần khởi động đầu, gửi lượt chấm và nhịp sống, nhận
bản cập nhật. Nhân viên phải đồng ý xử lý dữ liệu khuôn mặt trước khi được gán vào máy, và rút đồng ý thì mẫu bị xóa trên mọi máy.

\begin{figure}[!htb]
\centering
\begin{tikzpicture}[
  font=\footnotesize,
  uc/.style={draw, ellipse, align=center, minimum width=30mm, minimum height=8.5mm, inner sep=0pt}]
\draw (-3.65,-4.3) rectangle (5.5,4.35);
\node[anchor=north west] at (-3.6,4.32) {Máy chấm công và quản lý máy};
\node[uc] (c1) at (-2.0,3.45) {Chấm công tại máy};
\node[uc] (c3) at (-2.0,1.05) {Đăng ký khuôn mặt\\tại máy};
\node[uc] (c2) at (-2.0,2.25) {Xác thực khuôn mặt};
\node[uc] (c5) at (-2.0,-0.8) {Đồng ý hoặc rút\\đồng ý dữ liệu mặt};
\node[uc] (c4) at (-2.0,-2.9) {Gán người vào máy};
\node[uc] (d3) at (3.85,3.55) {Theo dõi trạng thái};
\node[uc] (d4) at (3.85,2.4) {Gửi lệnh tới máy};
\node[uc] (d5) at (3.85,1.25) {Đồng bộ lại\\danh sách người};
\node[uc] (d2) at (3.85,0.1) {Thu hồi máy};
\node[uc] (d1) at (3.85,-1.05) {Duyệt máy mới};
\node[uc] (d6) at (3.85,-2.2) {Phát hành bản\\cập nhật OTA};
\node[uc] (d7) at (3.85,-3.45) {Gửi lượt chấm,\\nhịp sống};
\actor{-5.4,2.0}{e}{Nhân viên}
\actor{-5.4,-3.0}{h}{Nhân sự}
\actor{7.0,1.6}{q}{Quản trị}
\actor{7.0,-3.0}{k}{Máy chấm\\công}
\draw (e) -- (c1.west);
\draw (e) -- (c3.west);
\draw (e) -- (c5.west);
\draw (h) -- (c4.west);
\draw[incl] (c1) -- node[font=\scriptsize, right] {$\ll$include$\gg$} (c2);
\draw[incl] (c3) -- node[font=\scriptsize, right] {$\ll$include$\gg$} (c2);
\foreach \u in {d3,d4,d5,d2,d1,d6} {\draw (q) -- (\u.east);}
\foreach \u in {d1,d6,d7} {\draw (k) -- (\u.east);}
\end{tikzpicture}
\caption{Phân rã ca sử dụng của máy chấm công và quản lý máy.}
\label{fig:uc_may}
\end{figure}

Hình~\ref{fig:uc_ban_giay} là phần của các vai bàn giấy. Mỗi loại việc chờ duyệt có đúng một bàn: đề nghị sửa hồ sơ, lượt chấm bị
giữ và tạm ứng do nhân sự duyệt, khiếu nại phiếu lương và người phụ thuộc do kế toán lương trả lời, còn giấy xác nhận thì cả hai bàn
cùng cấp. Quản lý chỉ duyệt đơn và xem công của người trong cây dưới quyền.

\begin{figure}[!htb]
\centering
\begin{tikzpicture}[
  font=\footnotesize,
  uc/.style={draw, ellipse, align=center, minimum width=31mm, minimum height=8.5mm, inner sep=0pt}]
\draw (-3.5,-6.6) rectangle (5.6,4.85);
\node[anchor=north west] at (-3.45,4.82) {Nhân sự, chấm công và lương};
\node[uc] (l1) at (-1.75,3.85) {Duyệt đơn};
\node[uc] (l2) at (-1.75,2.8) {Xem công của đội};
\node[uc] (l3) at (-1.75,1.75) {Quản lý hồ sơ, nhận\\việc, nghỉ việc};
\node[uc] (l4) at (-1.75,0.7) {Đổi mức lương};
\node[uc] (l5) at (-1.75,-0.35) {Duyệt đề nghị\\sửa hồ sơ};
\node[uc] (l6) at (-1.75,-1.4) {Khai ca, ngày lễ,\\loại phép};
\node[uc] (l7) at (-1.75,-2.45) {Xem và sửa\\bảng công};
\node[uc] (l8) at (-1.75,-3.5) {Duyệt lượt chấm\\bị giữ};
\node[uc] (l9) at (-1.75,-4.55) {Duyệt tạm ứng};
\node[uc] (l10) at (-1.75,-5.6) {Phát hành tài liệu\\cần ký};
\node[uc] (l11) at (1.95,-5.6) {Cấp giấy xác nhận};
\node[uc] (r1) at (3.85,3.85) {Khai chính sách\\lương, phụ cấp};
\node[uc] (r2) at (3.85,2.8) {Chạy và chốt\\kỳ lương};
\node[uc] (r3) at (3.85,1.75) {Chi tạm ứng};
\node[uc] (r4) at (3.85,0.7) {Trả lời khiếu nại\\phiếu lương};
\node[uc] (r5) at (3.85,-0.35) {Duyệt người\\phụ thuộc};
\node[uc] (r6) at (3.85,-1.4) {Xuất báo cáo\\bảo hiểm, thuế};
\actor{-5.6,3.2}{m}{Quản lý}
\actor{-5.6,-1.6}{h}{Nhân sự}
\actor{7.3,-3.6}{p}{Kế toán lương}
\draw (m) -- (l1.west);
\draw (m) -- (l2.west);
\foreach \u in {l1,l3,l4,l5,l6,l7,l8,l9,l10,l11} {\draw (h) -- (\u.west);}
\foreach \u in {r1,r2,r3,r4,r5,r6} {\draw (p) -- (\u.east);}
\draw (p) -- (l11.east);
\end{tikzpicture}
\caption{Phân rã ca sử dụng của các vai bàn giấy.}
\label{fig:uc_ban_giay}
\end{figure}

Hình~\ref{fig:uc_nhan_vien} là phần nhân viên làm trên điện thoại. Quản lý cũng là nhân viên nên có đủ các ca này. Năm loại đơn đi
chung một ca vì chúng đi cùng một luồng gửi và duyệt.

\begin{figure}[!htb]
\centering
\begin{tikzpicture}[
  font=\footnotesize,
  uc/.style={draw, ellipse, align=center, minimum width=52mm, minimum height=7.6mm, inner sep=0pt}]
\draw (-2.15,-10.35) rectangle (4.05,0.55);
\node[anchor=north west] at (-1.95,0.52) {Cổng nhân viên};
\foreach \y/\n/\t in {
  -0.4/u1/{Xem trang của tôi},
  -1.3/u2/{Xem công của tôi},
  -2.2/u3/{Xem lịch ca},
  -3.1/u4/{Gửi đơn: nghỉ, làm thêm,\\giải trình, công tác, từ xa},
  -4.0/u5/{Xin tạm ứng lương},
  -4.9/u6/{Xin giấy xác nhận},
  -5.8/u7/{Đề nghị sửa hồ sơ},
  -6.7/u8/{Xem phiếu lương},
  -7.6/u9/{Khiếu nại phiếu lương},
  -8.5/u10/{Đọc và ký tài liệu},
  -9.4/u11/{Nhận thông báo đẩy}} {
  \node[uc] (\n) at (0.95,\y) {\t};
}
\actor{-4.6,-3.4}{e}{Nhân viên}
\actor{-4.6,-7.6}{m}{Quản lý}
\foreach \u in {u1,u2,u3,u4,u5,u6,u7,u8,u9,u10,u11} {\draw (e) -- (\u.west);}
\draw[-{Triangle[open, length=2.6mm, width=2.4mm]}] (-4.6,-7.25) -- (-4.6,-4.4);
\end{tikzpicture}
\caption{Phân rã ca sử dụng của nhân viên trên điện thoại.}
\label{fig:uc_nhan_vien}
\end{figure}

Năm bảng dưới đặc tả các ca có nhiều bên tham gia nhất. Bảng~\ref{tab:uc_cham} là ca
chính của máy chấm công. Máy ghi lượt xuống flash trước khi gửi, nên ca này không phụ
thuộc mạng (F5).

\begin{table}[!htb]
\centering
\small
\caption{Đặc tả ca sử dụng Chấm công tại máy.}
\label{tab:uc_cham}
\begin{tabular}{|L{2.8cm}|L{10.8cm}|}
\hline
\textbf{Tác nhân} & Nhân viên \\
\hline
\textbf{Tiền điều kiện} & Người đó đã có mẫu khuôn mặt trên máy, máy đã nạp được gói mô hình \\
\hline
\textbf{Luồng chính} & (1)~Người bước vào tầm cảm biến, máy thức. (2)~Máy phát hiện mặt trong khung ngắm và chờ hộp mặt ổn định qua 2 lần phát hiện. (3)~Máy chấm mặt là thật. (4)~Máy tạo vector đặc trưng, so với các mẫu, điểm cao nhất từ 0,45 trở lên. (5)~Máy mở thanh chắn, kêu bíp, hiện tên. (6)~Máy ghi lượt chấm xuống flash rồi gửi lên máy chủ \\
\hline
\textbf{Luồng rẽ} & Mặt giả: sau 3 lần thử, máy báo từ chối. Không nhận ra: sau 2 lần thử, máy báo chưa nhận ra. Người vừa chấm rời đi rồi quay lại trong 5 phút: thanh chắn vẫn mở, không sinh lượt mới. Mất mạng: lượt nằm trong flash, có mạng thì gửi \\
\hline
\textbf{Hậu điều kiện} & Có một lượt chấm trên flash, sau đó có đúng một dòng trên máy chủ \\
\hline
\end{tabular}
\end{table}

Bảng~\ref{tab:uc_dang_ky} là ca nối máy với máy chủ qua khuôn mặt. Người vận hành chọn tên
trên màn hình, không gõ mã, vì gõ nhầm là buộc khuôn mặt người này vào hồ sơ người kia và
lỗi đó chỉ lộ ra khi đối chiếu bảng lương.

\begin{table}[!htb]
\centering
\small
\caption{Đặc tả ca sử dụng Đăng ký khuôn mặt tại máy.}
\label{tab:uc_dang_ky}
\begin{tabular}{|L{2.8cm}|L{10.8cm}|}
\hline
\textbf{Tác nhân} & Nhân sự, nhân viên, máy chấm công \\
\hline
\textbf{Tiền điều kiện} & Nhân viên đã đồng ý cho xử lý dữ liệu sinh trắc \\
\hline
\textbf{Luồng chính} & (1)~Nhân sự gán người vào máy trên dashboard, máy chủ đẩy tên xuống danh sách chờ của máy. (2)~Trên máy, người vận hành chọn tên ở màn Thêm người. (3)~Người được chụp nhìn thẳng, nghiêng trái, nghiêng phải; mỗi mẫu giữ được máy kêu một tiếng bíp. (4)~Đủ 3 mẫu, máy ghi bảng mẫu xuống flash. (5)~Máy gửi từng mẫu lên máy chủ tới khi broker xác nhận. (6)~Máy chủ mã hóa mẫu và đánh dấu cặp máy, người đã đăng ký \\
\hline
\textbf{Luồng rẽ} & Chưa có đồng ý: máy chủ từ chối gán với mã \texttt{BIOMETRIC\_CONSENT\_MISSING}. Bỏ dở: các mẫu đã lấy bị xóa. Mất mạng: mẫu mang cờ chưa báo và được gửi khi nối lại \\
\hline
\textbf{Hậu điều kiện} & Máy nhận ra người đó, máy chủ biết người đó đã có mặt trên máy nào \\
\hline
\end{tabular}
\end{table}

Bảng~\ref{tab:uc_don} là ca đơn nghỉ. Số dư bị giữ chỗ ngay lúc gửi, nên hai đơn cùng chờ
duyệt không tiêu quá số ngày còn lại.

\begin{table}[!htb]
\centering
\small
\caption{Đặc tả ca sử dụng Gửi và duyệt đơn nghỉ.}
\label{tab:uc_don}
\begin{tabular}{|L{2.8cm}|L{10.8cm}|}
\hline
\textbf{Tác nhân} & Nhân viên, quản lý hoặc nhân sự \\
\hline
\textbf{Tiền điều kiện} & Đã khai loại phép, người gửi có hồ sơ \\
\hline
\textbf{Luồng chính} & (1)~Nhân viên chọn loại phép và ngày. (2)~Biểu mẫu hiện số ngày làm việc bị trừ và số dư sau đơn do máy chủ tính. (3)~Nhân viên gửi đơn, số dư giữ chỗ trong \texttt{pending}. (4)~Người duyệt thấy đơn ở hộp chờ duyệt kèm số dư và người đang nghỉ. (5)~Người duyệt chấp nhận, số ngày chuyển sang \texttt{taken} và ghi vào ngày công \\
\hline
\textbf{Luồng rẽ} & Thiếu số dư: mã \texttt{LEAVE\_BALANCE\_SHORT}. Chồng ngày: mã \texttt{LEAVE\_OVERLAP}. Tự duyệt đơn của mình: mã \texttt{SELF\_DECISION}. Từ chối hoặc hủy: trả lại phần giữ chỗ \\
\hline
\textbf{Hậu điều kiện} & Số dư và ngày công thay đổi trong cùng một giao dịch \\
\hline
\end{tabular}
\end{table}

Hình~\ref{fig:hd_don} là quy trình của ca này. Người duyệt là quản lý trực tiếp hoặc người được quản lý ủy quyền, và đơn của người không có
ai phía trên thì tới bàn nhân sự. Ngày công đã dựng mà rơi vào đơn vừa duyệt được đổi ngay trong giao dịch duyệt, nên bảng công không
phải chờ lượt dựng đêm sau. Nhân viên hủy đơn khi chưa ai duyệt thì phần giữ chỗ cũng được trả lại.

\begin{figure}[!htb]
\centering
\begin{tikzpicture}[
  font=\footnotesize,
  act/.style={draw, rounded corners=2mm, align=center, text width=40mm, minimum height=8mm, inner sep=2pt},
  dec/.style={draw, diamond, aspect=1.6, inner sep=1pt, minimum width=7mm},
  arr/.style={-{Latex[length=2mm]}},
  lab/.style={font=\scriptsize, fill=white, inner sep=1pt}]
\foreach \x/\t in {0/Nhân viên, 4.8/Máy chủ, 9.6/Người duyệt} {
  \node at (\x,0.55) {\textbf{\t}};
}
\draw (-2.4,0.95) rectangle (12.0,-10.2);
\draw (-2.4,0.2) -- (12.0,0.2);
\draw[densely dashed] (2.4,0.95) -- (2.4,-10.2);
\draw[densely dashed] (7.2,0.95) -- (7.2,-10.2);
\fill (0,-0.3) circle (1.3mm);
\node[act] (a1) at (0,-1.2) {Chọn loại phép và ngày};
\node[act] (a2) at (4.8,-1.2) {Tính số ngày làm việc\\và số dư sau đơn};
\node[act] (a3) at (0,-2.5) {Gửi đơn};
\node[act] (a4) at (4.8,-2.5) {Kiểm số dư và\\đơn chồng ngày};
\node[dec] (d1) at (4.8,-3.65) {};
\node[act] (a5) at (0,-3.65) {Nhận mã lỗi};
\node[act] (a6) at (4.8,-5.0) {Giữ chỗ số dư, gửi đơn\\tới quản lý, người được\\ủy quyền hoặc nhân sự};
\node[act] (a7) at (9.6,-5.0) {Xem đơn kèm số dư\\và người đang nghỉ};
\node[dec] (d2) at (9.6,-6.35) {};
\node[act] (a8) at (4.8,-7.2) {Chuyển giữ chỗ sang\\đã dùng, đổi ngày\\đã dựng};
\node[act] (a9) at (4.8,-8.5) {Trả lại phần giữ chỗ};
\node[dec] (m1) at (0,-7.85) {};
\node[act] (a11) at (0,-8.9) {Nhận thông báo\\kết quả};
\draw (0,-4.55) circle (1.3mm); \fill (0,-4.55) circle (0.8mm);
\draw (0,-9.8) circle (1.3mm); \fill (0,-9.8) circle (0.8mm);
\draw[arr] (0,-0.43) -- (a1);
\draw[arr] (a1) -- (a2);
\draw[arr] (a2.south west) -- (a3.north east);
\draw[arr] (a3) -- (a4);
\draw[arr] (a4) -- (d1);
\draw[arr] (d1) -- node[lab, above] {[không đạt]} (a5);
\draw[arr] (a5) -- (0,-4.42);
\draw[arr] (d1) -- node[lab, right, pos=0.4] {[đạt]} (a6);
\draw[arr] (a6) -- (a7);
\draw[arr] (a7) -- (d2);
\draw[arr] (d2) |- node[lab, pos=0.3, right] {[duyệt]} (a8);
\draw[arr] (d2.east) -- ++(0.75,0) |- node[lab, pos=0.25, right] {[từ chối]} (a9);
\draw[arr] (a8.west) -| (m1.north);
\draw[arr] (a9.west) -- ++(-0.45,0) |- (m1.east);
\draw[arr] (m1) -- (a11);
\draw[arr] (a11) -- (0,-9.67);
\end{tikzpicture}
\caption{Quy trình gửi và duyệt đơn nghỉ.}
\label{fig:hd_don}
\end{figure}

Hình~\ref{fig:td_duyet_don} là bước duyệt của ca này, vẽ theo ba loại lớp phân tích. Lớp biên (boundary) là chỗ tác nhân thao tác,
tức hộp chờ duyệt trên dashboard. Lớp điều khiển (control) giữ luật của ca, gồm \texttt{LeaveController} nhận yêu cầu và
\texttt{LeaveService} kiểm quyền rồi ghi. Lớp thực thể (entity) là các bảng mà phép duyệt ghi vào. Lệnh đổi trạng thái chỉ khớp khi đơn
còn \texttt{PENDING}, nên hai người cùng bấm duyệt thì người sau nhận mã \texttt{REQUEST\_ALREADY\_DECIDED} và số dư không bị trừ hai
lần. Lệnh này cùng phép trừ số dư và phép đổi ngày công chạy trong một giao dịch.

\begin{figure}[!htb]
\centering
\begin{tikzpicture}[font=\footnotesize]
\sqactor{0}{Người duyệt}{-11.4}
\bceboundary{2.05}{Hộp chờ\\duyệt}{-11.4}
\bcecontrol{4.1}{\texttt{Leave}\\\texttt{Controller}}{-11.4}
\bcecontrol{6.15}{\texttt{Leave}\\\texttt{Service}}{-11.4}
\bceentity{8.2}{\texttt{Request}}{-11.4}
\bceentity{10.25}{\texttt{Leave}\\\texttt{Balance}}{-11.4}
\bceentity{12.3}{\texttt{Attendance}\\\texttt{Day}}{-11.4}
\sqexec{2.05}{-1.3}{-11.1}
\sqexec{4.1}{-2.3}{-10.6}
\sqexec{6.15}{-2.9}{-10.1}
\draw[sqsync] (0,-1.3) -- node[sqlab] {bấm Duyệt} (1.97,-1.3);
\draw[sqsync] (2.13,-2.3) -- node[sqlab] {\texttt{POST}\\\texttt{/requests/}\\\texttt{:id/decide}} (4.02,-2.3);
\draw[sqsync] (4.18,-2.9) -- node[sqlab] {\texttt{decide()}} (6.07,-2.9);
\sqexec{8.2}{-3.5}{-3.95}
\draw[sqsync] (6.23,-3.5) -- node[sqlab] {đọc đơn} (8.12,-3.5);
\draw[sqreply] (8.12,-3.95) -- node[sqlab] {đơn} (6.23,-3.95);
\sqselfr[0]{6.15}{-4.35}{kiểm quyền}
\sqexec{8.2}{-5.8}{-6.25}
\draw[sqsync] (6.23,-5.8) -- node[sqlab] {đổi trạng thái\\nếu \texttt{PENDING}} (8.12,-5.8);
\draw[sqreply] (8.12,-6.25) -- node[sqlab] {1 dòng} (6.23,-6.25);
\sqfrag{5.45}{-6.6}{13.1}{-9.75}{alt}
\node[sqnote, anchor=west] at (6.3,-6.85) {[duyệt]};
\sqexec{10.25}{-7.45}{-7.8}
\draw[sqsync] (6.23,-7.45) -- node[sqlab, pos=0.75] {\texttt{pending} giảm,\\\texttt{taken} tăng} (10.17,-7.45);
\sqexec{12.3}{-8.1}{-8.45}
\draw[sqsync] (6.23,-8.1) -- node[sqlab, pos=0.86] {ngày vắng\\thành nghỉ} (12.22,-8.1);
\draw[densely dashed] (5.45,-8.7) -- (13.1,-8.7);
\node[sqnote, anchor=west] at (6.3,-8.9) {[từ chối]};
\sqexec{10.25}{-9.25}{-9.55}
\draw[sqsync] (6.23,-9.25) -- node[sqlab, pos=0.75] {\texttt{pending} giảm} (10.17,-9.25);
\draw[sqreply] (6.07,-10.1) -- node[sqlab] {đơn đã quyết} (4.18,-10.1);
\draw[sqreply] (4.02,-10.6) -- node[sqlab] {201} (2.13,-10.6);
\draw[sqreply] (1.97,-11.1) -- node[sqlab] {báo đã duyệt} (0,-11.1);
\end{tikzpicture}
\caption{Trình tự duyệt một đơn nghỉ theo lớp biên, điều khiển và thực thể.}
\label{fig:td_duyet_don}
\end{figure}

Bảng~\ref{tab:uc_luong} là ca kỳ lương. Mỗi kỳ thuộc đúng một pháp nhân, vì lương, bảo hiểm,
thuế và chuyển khoản đều đi theo bên đứng tên hợp đồng. Một kỳ chạy được nhiều lượt nháp,
nhưng chỉ khi chốt thì phiếu mới là bản chính thức.

\begin{table}[!htb]
\centering
\small
\caption{Đặc tả ca sử dụng Chạy và chốt kỳ lương.}
\label{tab:uc_luong}
\begin{tabular}{|L{2.8cm}|L{10.8cm}|}
\hline
\textbf{Tác nhân} & Kế toán lương \\
\hline
\textbf{Tiền điều kiện} & Tháng lương đã bắt đầu, pháp nhân có chính sách hiệu lực \\
\hline
\textbf{Luồng chính} & (1)~Mở kỳ cho một pháp nhân. (2)~Chạy một lượt tính, việc nền sinh phiếu nháp. (3)~Soát phiếu, sửa đầu vào rồi chạy lại nếu cần. (4)~Sau khi hết tháng, chốt kỳ: đầu vào đóng băng, phiếu mang bản chép tài khoản ngân hàng. (5)~Xuất tệp ngân hàng, đánh dấu đã trả, gửi phiếu cho nhân viên \\
\hline
\textbf{Luồng rẽ} & Thiếu chính sách: lượt không chạy và báo mã lỗi kèm đường sửa. Tháng chưa hết: mã \texttt{PERIOD\_NOT\_ENDED}. Lượt chấm tới sau khi chốt: chờ nhân sự duyệt và chỉ sửa bảng công; khoản chênh, nếu có, đi qua khiếu nại lương và vào kỳ sau như một khoản truy lĩnh \\
\hline
\textbf{Hậu điều kiện} & Phiếu lương chính thức không đổi nữa, mỗi người mỗi tháng đúng một phiếu thường \\
\hline
\end{tabular}
\end{table}

Hình~\ref{fig:hd_luong} là quy trình kỳ lương. Vòng soát và chạy lại lặp bao nhiêu lần cũng được vì phiếu còn là nháp, và chỉ lúc chốt
thì đầu vào mới bị đóng băng. Sau khi chốt, sai lệch không sửa vào phiếu đã phát mà đi qua khiếu nại và thành một khoản của kỳ sau.

\begin{figure}[!htb]
\centering
\begin{tikzpicture}[
  font=\footnotesize,
  act/.style={draw, rounded corners=2mm, align=center, text width=40mm, minimum height=8mm, inner sep=2pt},
  dec/.style={draw, diamond, aspect=1.6, inner sep=1pt, minimum width=7mm},
  arr/.style={-{Latex[length=2mm]}},
  lab/.style={font=\scriptsize, fill=white, inner sep=1pt}]
\foreach \x/\t in {0/Kế toán lương, 4.8/Máy chủ, 9.6/Nhân viên} {
  \node at (\x,0.55) {\textbf{\t}};
}
\draw (-2.7,0.95) rectangle (12.0,-11.15);
\draw (-2.7,0.2) -- (12.0,0.2);
\draw[densely dashed] (2.4,0.95) -- (2.4,-11.15);
\draw[densely dashed] (7.2,0.95) -- (7.2,-11.15);
\fill (0,-0.3) circle (1.3mm);
\node[act] (b1) at (0,-1.15) {Mở kỳ cho một\\pháp nhân};
\node[dec] (m1) at (0,-2.12) {};
\node[act] (b2) at (0,-3.0) {Chạy lượt tính};
\node[act] (b3) at (4.8,-3.0) {Tính theo lô 500 người,\\sinh phiếu nháp};
\node[act] (b4) at (0,-4.25) {Soát phiếu và danh\\mục cần xử lý};
\node[dec] (d1) at (0,-5.35) {};
\node[act] (b5) at (0,-6.45) {Chốt kỳ sau\\khi hết tháng};
\node[act] (b6) at (4.8,-6.45) {Đóng băng đầu vào,\\phát phiếu, chép tài\\khoản ngân hàng};
\node[act] (b7) at (0,-7.8) {Xuất tệp ngân hàng,\\đánh dấu đã trả};
\node[act] (b8) at (4.8,-7.8) {Gửi phiếu, báo\\cho từng người};
\node[act] (b9) at (9.6,-7.8) {Xem phiếu lương};
\node[dec] (d2) at (9.6,-8.9) {};
\node[act] (b10) at (9.6,-10.0) {Gửi khiếu nại};
\node[act] (b11) at (0,-10.0) {Trả lời, hạn 5 ngày};
\node[act] (b12) at (4.8,-10.0) {Ghi khoản chênh\\vào kỳ sau};
\draw (11.6,-8.9) circle (1.3mm); \fill (11.6,-8.9) circle (0.8mm);
\draw (7.0,-10.0) circle (1.3mm); \fill (7.0,-10.0) circle (0.8mm);
\draw[arr] (0,-0.43) -- (b1);
\draw[arr] (b1) -- (m1);
\draw[arr] (m1) -- (b2);
\draw[arr] (b2) -- (b3);
\draw[arr] (b3.south west) -- (b4.north east);
\draw[arr] (b4) -- (d1);
\draw[arr] (d1) -- node[lab, right] {[ổn]} (b5);
\draw[arr] (d1.west) -- node[lab, above] {[cần sửa]} ++(-2.0,0) |- (m1.west);
\draw[arr] (b5) -- (b6);
\draw[arr] (b6.south west) -- (b7.north east);
\draw[arr] (b7) -- (b8);
\draw[arr] (b8) -- (b9);
\draw[arr] (b9) -- (d2);
\draw[arr] (d2) -- node[lab, above] {[đúng]} (11.47,-8.9);
\draw[arr] (d2) -- node[lab, right] {[sai]} (b10);
\draw[arr] (b10.south) -- ++(0,-0.4) -| (b11.south);
\draw[arr] (b11) -- (b12);
\draw[arr] (b12) -- (6.87,-10.0);
\end{tikzpicture}
\caption{Quy trình chạy, chốt kỳ lương và khiếu nại.}
\label{fig:hd_luong}
\end{figure}

Bảng~\ref{tab:uc_duyet_may} là ca đưa một máy mới vào hệ thống. Mọi máy nhận đúng cùng một
ảnh firmware, nên token lô không phân biệt được máy thật với máy giả. Mã nhận máy chỉ hiện trên màn hình thật, buộc người duyệt đứng trước đúng cái máy đó.

\begin{table}[!htb]
\centering
\small
\caption{Đặc tả ca sử dụng Duyệt máy mới.}
\label{tab:uc_duyet_may}
\begin{tabular}{|L{2.8cm}|L{10.8cm}|}
\hline
\textbf{Tác nhân} & Quản trị, máy chấm công \\
\hline
\textbf{Tiền điều kiện} & Máy có Wi-Fi và chạy firmware mang token của lô \\
\hline
\textbf{Luồng chính} & (1)~Máy gọi \texttt{POST /devices/register} kèm token lô, \texttt{deviceId} và mã nhận máy 6 số, đồng thời hiện mã trên màn hình. (2)~Máy chủ tạo dòng \texttt{PENDING}. (3)~Quản trị đọc mã trên màn hình, gõ vào dashboard, đặt tên và vị trí. (4)~Lần hỏi kế tiếp, máy nhận JWT 90 ngày và nối broker \\
\hline
\textbf{Luồng rẽ} & Sai mã: không duyệt. Trong lúc chờ, máy vẫn chấm công và giữ lượt trong flash \\
\hline
\textbf{Hậu điều kiện} & Máy ở trạng thái \texttt{APPROVED}, có JWT riêng \\
\hline
\end{tabular}
\end{table}

\subsection{Yêu cầu chức năng}
\label{section:3.1.2}

Bảng~\ref{tab:yc_chuc_nang} liệt kê yêu cầu chức năng, tách từ bốn sản phẩm đã nêu ở
Mục~\ref{section:1.2}. F1--F6 thuộc máy chấm công, F7--F9 thuộc phần quản lý máy từ xa,
F10--F13 thuộc hệ thống nhân sự, F14 và F15 thuộc cổng nhân viên.

\begin{table}[!htb]
\centering
\small
\caption{Yêu cầu chức năng.}
\label{tab:yc_chuc_nang}
\begin{tabular}{|c|L{3.2cm}|L{9.6cm}|}
\hline
\textbf{Mã} & \textbf{Nhóm} & \textbf{Yêu cầu} \\
\hline
F1 & Máy chấm công & Chấm công bằng khuôn mặt không tiếp xúc, xử lý trên chip \\
\hline
F2 & Máy chấm công & Chặn ảnh in và ảnh hoặc video phát trên màn hình \\
\hline
F3 & Máy chấm công & Đăng ký khuôn mặt tại máy theo danh sách máy chủ giao \\
\hline
F4 & Máy chấm công & Tự thức khi có người, tự nghỉ khi không có ai \\
\hline
F5 & Máy chấm công & Vẫn chấm công khi mất mạng, tự gửi bù khi nối lại \\
\hline
F6 & Máy chấm công & Mở thanh chắn và phát tiếng báo \\
\hline
F7 & Quản lý máy & Duyệt máy mới bằng mã hiện trên màn hình, thu hồi máy \\
\hline
F8 & Quản lý máy & Theo dõi trạng thái, ra lệnh từ xa \\
\hline
F9 & Quản lý máy & Cập nhật firmware và mô hình qua OTA, tự quay về bản cũ khi hỏng \\
\hline
F10 & Nhân sự & Hồ sơ nhân viên, cơ cấu tổ chức, vai và phạm vi nhìn thấy \\
\hline
F11 & Nhân sự & Ca làm, ngày công, bảng công, sửa tay có vết \\
\hline
F12 & Nhân sự & Đơn nghỉ, đơn làm thêm, luồng duyệt, số dư phép \\
\hline
F13 & Nhân sự & Kỳ lương, tính thuế và bảo hiểm, phiếu lương, báo cáo \\
\hline
F14 & Cổng nhân viên & Xem công, phép, phiếu lương và gửi đơn trên điện thoại \\
\hline
F15 & Cổng nhân viên & Nhận thông báo đẩy, xem lại dữ liệu đã tải khi mất mạng \\
\hline
\end{tabular}
\end{table}

\subsection{Yêu cầu phi chức năng và ràng buộc}
\label{section:3.1.3}

Bảng~\ref{tab:yc_phi_chuc_nang} đặt các ràng buộc kiểm tra được. Phần lớn sinh ra từ
phần cứng và từ chỗ người chấm công chỉ đứng trước máy 1--2~s. Kết quả đo nằm ở
Chương~\ref{chapter:Experiment}. NF9 chưa có con số, vì đồ án chưa lập bảng giá linh
kiện.

\begin{table}[!htb]
\centering
\small
\caption{Yêu cầu phi chức năng.}
\label{tab:yc_phi_chuc_nang}
\begin{tabular}{|c|L{2.6cm}|L{10.2cm}|}
\hline
\textbf{Mã} & \textbf{Nhóm} & \textbf{Mục tiêu} \\
\hline
NF1 & Độ trễ & Một lượt chạy ba mô hình không quá 360~ms \\
\hline
NF2 & Bộ nhớ & Đáy SRAM trống từ 24~KB trở lên khi chạy đủ tính năng; ba mô hình vừa khe 2,25~MB \\
\hline
NF3 & Độ chính xác & Phát hiện: AP từ 0,90 trở lên trên mặt từ 38~px. Lượng tử hóa không làm thước chính của mỗi nhánh sụt quá 1\%. Chống giả mạo chặn mọi khung giả của tập OV5640 \\
\hline
NF4 & Riêng tư & Ảnh khuôn mặt không rời máy chấm công; vector đặc trưng được mã hóa khi lưu trên máy chủ \\
\hline
NF5 & Bền dữ liệu & Không mất lượt chấm khi mất điện hay mất mạng; gửi lại nhiều lần cũng chỉ thành một bản ghi \\
\hline
NF6 & Bảo mật & Mã hóa đường truyền, mỗi máy một danh tính, phân quyền theo dòng dữ liệu, xác thực hai lớp cho vai bàn giấy \\
\hline
NF7 & Quy mô & Dùng được với 5.000 nhân viên: danh sách có tìm, lọc, phân trang, truy vấn không quét cả bảng \\
\hline
NF8 & Khả dụng & Tiếng Việt và tiếng Anh ở cả máy lẫn web; cổng nhân viên dùng được trên điện thoại \\
\hline
NF9 & Chi phí & Máy chỉ dùng một vi điều khiển và các module bán sẵn \\
\hline
\end{tabular}
\end{table}

\section{Kiến trúc tổng thể}
\label{section:3.2}

Hệ thống gồm ba phần như Hình~\ref{fig:kien_truc}. Máy chấm công tự quyết định tại
chỗ: nhận diện, mở thanh chắn và ghi lượt chấm mà không cần mạng (F1, F5). Máy chủ chạy
trên một VPS bằng Docker Compose, giữ danh tính và nghiệp vụ. Người dùng làm việc qua
ứng dụng Next.js triển khai trên Vercel, và trình duyệt gọi thẳng API qua REST và
WebSocket. Mọi lưu lượng vào VPS đi qua Traefik, gồm cả cổng 8883 của MQTT. Máy nối
vào EMQX, EMQX hỏi API xem JWT của máy có hợp lệ không, và API nhận bản tin qua một tài
khoản dịch vụ của broker.

\begin{figure}[!htb]
\centering
\begin{tikzpicture}[
  font=\footnotesize,
  box/.style={draw, align=center, minimum height=9mm, inner sep=2pt},
  lab/.style={font=\scriptsize, align=center},
  arr/.style={-{Latex[length=2mm]}},
  both/.style={{Latex[length=2mm]}-{Latex[length=2mm]}}]
\node[box, text width=26mm, minimum height=32mm] (k) at (0,0) {Máy chấm công\\ESP32-S3\\[1mm]ba mô hình,\\bảng mẫu,\\nhật ký lượt chấm};
\draw[dashed] (2.75,-2.45) rectangle (9.85,2.45);
\node[anchor=north west] at (2.8,2.42) {VPS, Docker Compose};
\node[box, text width=11mm, minimum height=38mm] (t) at (3.6,-0.1) {Traefik};
\node[box, text width=20mm] (e) at (6.0,1.05) {EMQX};
\node[box, text width=20mm] (a) at (6.0,-0.75) {API NestJS};
\node[box, text width=19mm] (p) at (8.7,0.35) {PostgreSQL};
\node[box, text width=19mm] (r) at (8.7,-0.95) {Redis,\\BullMQ};
\node[box, text width=23mm] (v) at (12.3,1.1) {Next.js\\trên Vercel};
\node[box, text width=23mm] (u) at (12.3,-1.1) {Trình duyệt,\\PWA, Android};
\draw[arr] ([yshift=7mm]k.east) -- node[lab, above] {MQTTS} ([yshift=8mm]t.west);
\draw[arr] ([yshift=-7mm]k.east) -- node[lab, below] {HTTPS} ([yshift=-6mm]t.west);
\draw[arr] ([yshift=11.5mm]t.east) -- (e.west);
\draw[arr] ([yshift=-6.5mm]t.east) -- (a.west);
\draw[both] (e) -- node[lab, right] {bản tin,\\kiểm JWT} (a);
\draw[arr] (a.east) -- (p.west);
\draw[arr] (a.east) -- (r.west);
\draw[arr] (u) -- node[lab, right] {tải trang} (v);
\draw[arr] (u.west) -- ++(-0.45,0) |- node[lab, pos=0.75, below] {REST, WebSocket} ([yshift=-16mm]t.east);
\end{tikzpicture}
\caption{Kiến trúc tổng thể của hệ thống.}
\label{fig:kien_truc}
\end{figure}

Quyết định lớn nhất của phần này là ai giữ thứ gì, tóm tắt ở Bảng~\ref{tab:noi_giu}.
Máy chủ giữ danh tính vì một người tồn tại trong công ty nhiều năm, còn mẫu khuôn mặt
gắn với một máy cụ thể và mất khi máy hỏng. Máy không tự đặt mã nhân viên. Mã tăng dần làm hai máy cùng sinh mã 1 cho hai người, còn mã ở dải riêng vượt kiểu \texttt{Int} của máy chủ, khiến PostgreSQL từ chối bản ghi trong khi máy vẫn thấy broker đã nhận. Máy vì thế chỉ đăng ký người máy chủ đã giao.

\begin{table}[!htb]
\centering
\small
\caption{Nơi giữ từng loại dữ liệu.}
\label{tab:noi_giu}
\begin{tabular}{|L{4.0cm}|L{3.9cm}|L{5.6cm}|}
\hline
\textbf{Dữ liệu} & \textbf{Nguồn gốc} & \textbf{Bản ở nơi khác} \\
\hline
Nhân viên, mã, tên & Máy chủ & Mã và tên trong bảng mẫu của máy \\
\hline
Ảnh khuôn mặt & Bộ nhớ máy, chỉ khi xử lý & Không lưu, không gửi đi (NF4) \\
\hline
Mẫu khuôn mặt & Máy chụp & Bản mã hóa AES-256-GCM trên máy chủ, đẩy xuống các máy khác \\
\hline
Lượt chấm & Flash của máy & Máy chủ, chống trùng theo cặp máy và mã lượt \\
\hline
Ngưỡng mô hình & Tiêu đề gói mô hình & Đổi cùng mô hình, không có bản sao \\
\hline
Ngưỡng nghiệp vụ & NVS của máy & Đổi được bằng lệnh \texttt{SET\_CONFIG} \\
\hline
Chính sách lương & Máy chủ & Phiếu lương ghi dòng chính sách đã dùng \\
\hline
\end{tabular}
\end{table}

Máy và máy chủ viết bằng hai ngôn ngữ nên dễ lệch hợp đồng. Thư mục
\mbox{\texttt{contracts/}} là nguồn duy nhất của định dạng bản tin, bảng topic và phiên bản
mô hình. Công cụ sinh mã phát ra kiểu C cho firmware và kiểu TypeScript cho backend,
frontend, và CI chặn khi mã sinh ra lệch nguồn. Mã nguồn đặt trong một kho chung gồm
\texttt{firmware}, \texttt{ml}, \texttt{backend}, \texttt{frontend}, \texttt{android},
\mbox{\texttt{contracts}}, \texttt{deploy} và \texttt{hardware}.

Hình~\ref{fig:module_be} chia 27 module của backend thành sáu nhóm nghiệp vụ và một phần nền. Phụ thuộc chính đi theo dòng dữ liệu: lượt
chấm của nhóm máy chấm công thành ngày công của nhóm thời gian, ngày công thành dòng của nhóm lương. Mọi nhóm phát sự kiện cho thông
báo và realtime. Không module nào tự đọc biến môi trường hay mở kết nối riêng tới cơ sở dữ liệu, cả hai đều đi qua phần nền.

\begin{figure}[!htb]
\centering
\begin{tikzpicture}[
  font=\footnotesize,
  grp/.style={draw, align=center, text width=40mm, minimum height=19mm, inner sep=2pt},
  ext/.style={draw, rounded corners=2mm, align=center, minimum height=8mm, inner sep=3pt},
  arr/.style={-{Latex[length=2mm]}},
  lab/.style={font=\scriptsize, fill=white, inner sep=1pt}]
\node[ext, text width=40mm] (web) at (2.2,3.25) {Trình duyệt, PWA, ứng dụng Android};
\node[ext, text width=40mm] (kio) at (6.8,3.25) {Máy chấm công qua EMQX};
\node[grp] (g1) at (2.2,1.2) {\textbf{Danh tính và quyền}\\[1mm]{\scriptsize\ttfamily auth, users, audit}};
\node[grp] (g2) at (6.8,1.2) {\textbf{Máy chấm công}\\[1mm]{\scriptsize\ttfamily devices, enrollment,\\models, mqtt, attendance}};
\node[grp] (g3) at (11.4,1.2) {\textbf{Thời gian}\\[1mm]{\scriptsize\ttfamily shifts, timesheet, leave}};
\node[grp] (g4) at (2.2,-1.25) {\textbf{Nhân sự}\\[1mm]{\scriptsize\ttfamily employees, org, onboarding,\\profile, documents,\\certificates, assets, search}};
\node[grp] (g5) at (6.8,-1.25) {\textbf{Lương}\\[1mm]{\scriptsize\ttfamily payroll, compensation,\\policy, disputes}};
\node[grp] (g6) at (11.4,-1.25) {\textbf{Thông báo, báo cáo}\\[1mm]{\scriptsize\ttfamily notifications, realtime,\\reports}};
\node[draw, align=center, text width=128mm, minimum height=10mm, inner sep=2pt] (base) at (6.8,-3.35) {\textbf{Nền dùng chung}: guard và vai, phạm vi dữ liệu, interceptor nhật ký và sự kiện, cache;\\Prisma; hàng đợi BullMQ; kiểm biến môi trường; {\scriptsize\ttfamily health}};
\node[ext, text width=24mm] (pg) at (4.6,-4.95) {PostgreSQL};
\node[ext, text width=24mm] (rd) at (9.0,-4.95) {Redis};
\draw[arr] (web) -- node[lab, right] {REST, WebSocket} (g1);
\draw[arr] (kio) -- node[lab, right] {bản tin MQTT} (g2);
\draw[arr] (g2) -- node[lab, above] {lượt chấm} (g3);
\draw[arr] (g3.south west) -- node[lab, right, pos=0.4] {ngày công} (g5.north east);
\draw[arr] (g4) -- node[lab, above] {hồ sơ} (g5);
\draw[arr] (g5) -- node[lab, above] {sự kiện} (g6);
\draw[arr] (base) -- (pg);
\draw[arr] (base) -- (rd);
\end{tikzpicture}
\caption{Các nhóm module của backend.}
\label{fig:module_be}
\end{figure}

\section{Thiết kế phần cứng}
\label{section:3.3}

\subsection{Sơ đồ khối và phân bổ bus}
\label{section:3.3.1}

Hình~\ref{fig:khoi_phan_cung} là sơ đồ khối của máy, linh kiện và lý do chọn đã nêu ở
Bảng~\ref{tab:ngoai_vi}. Board ESP32-S3-CAM mang sẵn module N16R8 và camera OV5640, các
module còn lại cắm lên một mạch đế. Camera chiếm 14 chân, màn hình 7 chân, và sau khi gán
hết ngoại vi thì board không còn chân GPIO thường nào. Ba đường ít đổi trạng thái là reset
cảm ứng, XSHUT của cảm biến khoảng cách và SD của bộ khuếch đại vì thế đi qua PCF8574.
Chân của PCF8574 lên mức cao lúc cấp nguồn, nên nó chỉ nhận những việc mà mức cao lúc khởi
động là vô hại: RST và XSHUT thả cao thì chip chạy, SD cao thì bộ khuếch đại thức nhưng
chưa có dữ liệu I\textsuperscript{2}S. Bảng chân đầy đủ nằm ở Phụ lục~\ref{appendix:chan}.

\begin{figure}[!htb]
\centering
\begin{tikzpicture}[
  font=\footnotesize,
  box/.style={draw, align=center, text width=28mm, minimum height=8mm, inner sep=2pt},
  lab/.style={font=\scriptsize, above},
  arr/.style={-{Latex[length=2mm]}}]
\node[draw, align=center, text width=31mm, minimum height=46mm] (c) at (0,0) {ESP32-S3-CAM\\[1mm]\mbox{WROOM-1 N16R8}\\flash 16~MB\\PSRAM 8~MB};
\node[box] (cam) at (-5.2,1.65) {OV5640};
\node[box] (lcd) at (-5.2,0.35) {ST7796S, 4 inch};
\node[box] (amp) at (-5.2,-0.95) {MAX98357A và loa};
\node[box] (srv) at (-5.2,-2.25) {Servo SG90};
\draw[arr] (cam.east) -- node[lab] {DVP, 14 chân} ([yshift=16.5mm]c.west);
\draw[arr] ([yshift=3.5mm]c.west) -- node[lab] {SPI 80~MHz} (lcd.east);
\draw[arr] ([yshift=-9.5mm]c.west) -- node[lab] {I\textsuperscript{2}S} (amp.east);
\draw[arr] ([yshift=-22.5mm]c.west) -- node[lab] {PWM 50~Hz} (srv.east);
\draw[thick] (3.0,2.1) -- (3.0,-2.4);
\draw (c.east) -- node[lab] {I\textsuperscript{2}C} (3.0,0);
\node[box] (gt) at (5.2,1.95) {GT911};
\node[box] (tof) at (5.2,0.65) {VL53L1X};
\node[box] (io) at (5.2,-0.65) {PCF8574:\\RST, XSHUT, SD};
\node[box] (rtc) at (5.2,-1.95) {DS3231};
\draw (3.0,1.95) -- (gt.west);
\draw (3.0,0.65) -- (tof.west);
\draw (3.0,-0.65) -- (io.west);
\draw (3.0,-1.95) -- (rtc.west);
\end{tikzpicture}
\caption{Sơ đồ khối phần cứng của máy chấm công.}
\label{fig:khoi_phan_cung}
\end{figure}

Bus I\textsuperscript{2}C chạy 400~kHz và do bốn thiết bị dùng chung. Mỗi đường có đúng một
điện trở kéo lên: SDA dùng R1 4,7~k$\Omega$ trên mạch đế, SCL dùng điện trở có sẵn của
module cảm ứng. GT911, VL53L1X và DS3231 đều đạt 400~kHz, nhưng datasheet của PCF8574 chỉ
ghi tối đa 100~kHz~\cite{pcf8574}. Bài kiểm tra bus đạt 7/7 ca ở 400~kHz
với đủ thiết bị, nên thiết kế giữ mức này. Đây là một điểm nằm ngoài đặc tả: kết quả đo
chỉ nói về linh kiện đang dùng, đổi lô PCF8574 thì phải chạy lại bài kiểm tra.

\subsection{Nguồn điện}
\label{section:3.3.2}

Máy dùng hai đường 5~V chung mass, tóm tắt ở Bảng~\ref{tab:nguon}, lấy từ hai adapter hoặc từ
khối pin ở Mục~\ref{section:4.2}. Hai đường tách nhau
vì hai đỉnh dòng trùng nhau: máy phát tiếng đúng lúc mở thanh chắn. Đặt chung một đường
thì đỉnh 1,34~A cộng 0,7~A của servo đủ kéo sụt áp và làm ESP32 khởi động lại. Đường 2
không được là sạc dự phòng, vì sạc tự ngắt khi tải nhỏ mà servo đứng yên gần như không
ăn dòng.

\begin{table}[!htb]
\centering
\small
\caption{Các đường nguồn của máy chấm công.}
\label{tab:nguon}
\begin{tabular}{|L{2.3cm}|L{2.2cm}|L{4.3cm}|L{2.0cm}|L{1.7cm}|}
\hline
\textbf{Đường} & \textbf{Mức} & \textbf{Tải} & \textbf{Điển hình} & \textbf{Đỉnh} \\
\hline
Đường 1 & 5~V, 2~A & ESP32-S3, camera, màn hình, khuếch đại, cảm biến & 375~mA & 1,34~A \\
\hline
Đường 2 & 5~V, từ 1~A & Servo SG90 & 150~mA & 700~mA \\
\hline
\texttt{+3V3\_LCD} & 3,3~V, từ đường 1 & Tấm màn hình & 100~mA & 150~mA \\
\hline
\end{tabular}
\end{table}

Tấm màn hình có đường 3,3~V riêng qua module ổn áp U2. Dãy LED đèn nền không tự ổn áp nên
độ sáng bám theo điện áp đặt vào, trong khi đường 3,3~V của board còn nuôi lõi 240~MHz và
PSRAM. Chung đường đó thì gợn điện áp theo tải của chip đi thẳng vào dãy LED, và màn hình nhấp nháy thấy rõ trên nền trắng. Đất của màn hình cũng đóng lại ngay tại chân đất U2. Các tụ lọc gồm 1000~$\mu$F gần đầu vào, 470~$\mu$F gần bộ khuếch đại, gần servo và
ở đầu vào U2, 100~$\mu$F ở đầu ra U2.

Hình~\ref{fig:nguon} vẽ cây nguồn. Board ESP32-S3-CAM và bộ khuếch đại nhận 5~V bằng hai nhánh
riêng từ J10. U2 là tải duy nhất được mượn nhánh của tải khác, vì nó ăn 100~mA đều đặn, không có
đỉnh để đẩy sang ai. Ba cảm biến I\textsuperscript{2}C và R1 vẫn ở đường 3,3~V của board.

\begin{figure}[!htb]
\centering
\begin{tikzpicture}[
  font=\footnotesize,
  box/.style={draw, align=center, text width=33mm, minimum height=9mm, inner sep=2pt},
  lab/.style={font=\scriptsize, above},
  arr/.style={-{Latex[length=2mm]}}]
\node[box] (j10) at (0,0.35) {J10: đường 1\\5~V, 2~A};
\node[box] (u1) at (4.8,1.75) {U1: \mbox{ESP32-S3-CAM}};
\node[box] (u2) at (4.8,0.35) {U2: ổn áp 3,3~V};
\node[box] (j7) at (4.8,-1.05) {J7: MAX98357A};
\node[box] (sen) at (10.1,1.75) {Cảm biến I\textsuperscript{2}C, R1\\(J4, J5, J6)};
\node[box] (j3) at (10.1,0.35) {J3: tấm màn hình};
\node[box] (j11) at (0,-2.3) {J11: đường 2\\5~V, từ 1~A};
\node[box] (j9) at (4.8,-2.3) {J9: SG90};
\draw[arr] (j10.east) -- (u1.west);
\draw[arr] (j10.east) -- (j7.west);
\draw[arr] (u1) -- (u2);
\draw[arr] (u1) -- node[lab] {\texttt{+3V3}} (sen);
\draw[arr] (u2) -- node[lab] {\texttt{+3V3\_LCD}} (j3);
\draw[arr] (j11) -- (j9);
\end{tikzpicture}
\caption{Cây nguồn của máy chấm công.}
\label{fig:nguon}
\end{figure}

\subsection{Sơ đồ nguyên lý và mạch in}
\label{section:3.3.3}

Mạch đế mới dừng ở bản thiết kế, bản đang chạy lắp trên breadboard (Mục~\ref{section:4.2}).
Hình~\ref{fig:nguyen_ly} là sơ đồ nguyên lý của mạch đế. U1 là board ESP32-S3-CAM, J3 là
module màn hình kèm cảm ứng, J4 là VL53L1X, J5 và J13 là hai hàng chân của PCF8574, J6 là
DS3231, J7 là MAX98357A, J9 là servo, còn J10 và J11 là hai đầu vào nguồn. J14 chừa sẵn lỗ
cho khe thẻ nhớ của màn hình, chưa nối dây. Mạch đế không có linh kiện dán: mọi chân đều
xuyên lỗ, các module cắm qua đế nên một module hỏng thì thay được ngay. Sơ đồ và mạch in
được sinh bằng script từ bảng chân, nên bảng chân là nguồn duy nhất.

\begin{figure}[!htb]
\centering
\includegraphics[width=\textwidth, trim=38 150 38 65, clip]{../../../hardware/export/kiosk_schematic_1eb1707.pdf}
\caption{Sơ đồ nguyên lý mạch đế.}
\label{fig:nguyen_ly}
\end{figure}

Hình~\ref{fig:mach_in} là mạch in. Mạch rộng 158~mm, cao 121~mm, hai lớp. Chiều cao đặt
theo tấm màn hình 108,04~mm cộng chỗ cho tên chân ở hai đầu. Vùng nét đứt là chỗ thân
module nằm: màn hình ở giữa, board ESP32-S3-CAM bên trái, hai đầu vào nguồn ở góc dưới bên
phải. Dây nguồn rộng 1,0~mm đi mặt sau, dây tín hiệu 0,25~mm đi mặt trước, tổng cộng 46
via. Via nguồn khoan 0,6~mm vì lỗ 0,4~mm chỉ chịu 0,87~A theo IPC-2221, dưới 0,91~A của via
mass nặng nhất. Bản hiện tại đạt 0 lỗi ERC và 0 vi phạm DRC.

\begin{figure}[!htb]
\centering
\includegraphics[width=0.9\textwidth, trim=189 142 190 141, clip]{../../../hardware/export/pcb_top.png}
\caption{Mạch đế, ảnh dựng từ KiCad.}
\label{fig:mach_in}
\end{figure}

Sơ đồ nguyên lý chỉ khai được rằng mọi chân GND cùng một nút, còn đồng thì có điện trở:
700~mA của servo chạy qua một đoạn đồng tạo hiệu điện thế trên đoạn đó, và MAX98357A khuếch
đại nhấp nhô ấy ra loa. Việc đi dây vì thế có luật riêng: (i)~hai mass gặp nhau tại đúng
một điểm cạnh hai đầu vào nguồn; (ii)~servo và bộ khuếch đại về thẳng đầu vào của mình;
(iii)~board ESP32-S3-CAM và bộ khuếch đại nhận 5~V bằng hai nhánh song song, vì nối chuỗi
thì đỉnh dòng của bộ khuếch đại chạy qua đoạn đồng nuôi chip; (iv)~tụ 470~$\mu$F đặt sát
chân tải.

\section{Thiết kế pipeline AI trên chip}
\label{section:3.4}

Pipeline xử lý từng khung theo Hình~\ref{fig:pipeline}, còn ba vùng ảnh của ba mô hình đã
vẽ ở Hình~\ref{fig:vung_anh}. Bộ phát hiện chạy trên mỗi khung AI lấy được, 32,9~ms một
lần. Hai mô hình sau tốn 73,3~ms và 246,4~ms nên chỉ chạy khi hộp mặt đã đứng yên qua 2 lần
phát hiện với IoU từ 0,5 trở lên (NF1). Kết quả của các cổng được báo ra ngay sau bước
phát hiện, trước hai mô hình chậm, qua một hàm quan sát. Nếu chỉ báo sau khi cả bước xong thì dòng Đang nhận diện gần như không bao giờ hiện, màn hình nhảy thẳng từ Lại gần hơn sang kết quả.

\begin{figure}[!htb]
\centering
\begin{tikzpicture}[
  font=\footnotesize,
  node distance=9mm and 3mm,
  box/.style={draw, align=center, text width=17mm, minimum height=12mm, inner sep=1.5pt},
  res/.style={draw, rounded corners, align=center, text width=18mm, minimum height=12mm, inner sep=1pt},
  arr/.style={-{Latex[length=2mm]}}]
\node[box] (f) {Khung mới};
\node[box, right=of f] (d) {YuNet};
\node[box, right=of d] (g) {Bốn cổng hình học};
\node[box, right=of g] (s) {Ổn định qua 2 lần};
\node[box, right=of s] (l) {Chống giả mạo};
\node[box, right=of l] (r) {Nhận diện, so mẫu};
\node[box, right=of r] (m) {Báo khớp};
\node[res, below=of g] (o1) {Nhắc đưa mặt vào khung, lại gần};
\node[res, below=of l] (o2) {3 lần thử thì báo giả mạo};
\node[res, below=of r] (o3) {2 lần thử thì báo chưa nhận ra};
\draw[arr] (f) -- (d);
\draw[arr] (d) -- (g);
\draw[arr] (g) -- (s);
\draw[arr] (s) -- (l);
\draw[arr] (l) -- (r);
\draw[arr] (r) -- (m);
\draw[arr] (g) -- node[font=\scriptsize, right] {không đạt} (o1);
\draw[arr] (l) -- node[font=\scriptsize, right] {dưới 0,50} (o2);
\draw[arr] (r) -- node[font=\scriptsize, right] {dưới 0,45} (o3);
\end{tikzpicture}
\caption{Luồng xử lý một khung của pipeline AI.}
\label{fig:pipeline}
\end{figure}

Bốn cổng hình học lọc mặt trước khi tốn thời gian cho hai mô hình sau: (i)~điểm phát hiện
từ 0,35; (ii)~ít nhất một nửa diện tích hộp mặt nằm trong khung ngắm; (iii)~mặt rộng từ
100~px trong khung; (iv)~ô vuông quanh mặt lọt trong khung. Cổng khung ngắm có mặt vì camera
chụp 480 cột còn màn hình chỉ hiện 213 cột giữa. Thiếu cổng này thì mặt nằm ngoài màn hình
vẫn giành được lượt của người đứng giữa khung, vì mặt mới được chọn theo cỡ lớn nhất. Mặt
đang bám chỉ tính là rời khung khi phần nằm trong khung xuống dưới 25\%, nên người đứng lấp
lửng ở mép không bị đóng mở lượt liên tục.

Máy bám một mặt chính. Mặt đang bám được giữ nếu lần phát hiện sau còn hộp trùng nó, nên
người thứ hai to hơn bước vào không cướp lượt của người đang được xác thực. Bộ phát hiện trượt lẻ tẻ ngay cả trên mặt đứng yên. Máy chỉ bỏ mặt đang bám sau 2 lần trượt liên tiếp.

Hai kết luận xấu có số lần thử khác nhau. Mặt thật thỉnh thoảng rơi dưới ngưỡng sống trên
một khung đơn lẻ, nên máy thử 3 lần mới báo giả mạo, mỗi lần thử cách nhau 3 lần phát hiện.
Không nhận ra thì thử 2 lần. Ba ngưỡng 0,35, 0,50 và 0,45 nằm trong tiêu đề gói mô hình nên
đổi cùng mô hình, và không đưa lên dashboard, vì hạ ngưỡng nhận diện là cho người lạ mở
thanh chắn. Các ngưỡng nghiệp vụ như cỡ mặt, tỉ lệ khung ngắm và thời gian chống trùng thì
nằm ở NVS.

Đăng ký đi qua chính đường xác thực này. Chống giả mạo chạy trước nhận diện, nên ảnh giơ lên
lúc đăng ký bị từ chối trước khi có vector đặc trưng và không mẫu giả nào vào được bảng. Máy
đo góc quay bằng năm điểm mốc: chiếu vector từ trung điểm hai mắt tới mũi lên trục hai mắt.
Mốc nhìn thẳng lấy trung bình của chính người đó, vì mốc thực lệch khoảng 0,06 so với 0, gần một phần ba ngưỡng quay 0,20. Nhờ phép đo này, người đứng yên
không thể lấy đủ ba mẫu chính diện gần trùng nhau.

Ba mô hình đóng chung một gói. Tiêu đề gói mang phiên bản, vị trí, sha256 và ba ngưỡng, và
gói nằm ở một trong hai phân vùng \texttt{models\_0}, \texttt{models\_1}. Hậu xử lý gồm giải
mã, NMS, căn chỉnh và chuẩn hóa viết bằng C++ trong \texttt{ai\_engine} và kiểm bằng bộ vector
vàng do Python sinh ở \texttt{contracts/golden/}, nên hai phía không lệch nhau.

Hình~\ref{fig:kien_truc_mo_hinh} là kiến trúc ba mô hình ở mức khối. Cả ba dựng trên tích chập depthwise và tích chập 1$\times$1, đúng
loại phép toán mà ESP-DL có kernel 8 bit. Nhận diện tốn nhiều phép tính nhất, 221,2 triệu MAC, vì nó giữ mặt ở 112$\times$112 tới bản đồ
7$\times$7 cuối cùng. Phát hiện chỉ tốn 25,2 triệu MAC trên ảnh 160$\times$120, và việc giải mã hộp, NMS chạy bằng C++ sau mô hình.

\begin{figure}[!htb]
\centering
\begin{tikzpicture}[
  font=\scriptsize,
  node distance=3.5mm,
  blk/.style={draw, align=center, text width=19.5mm, minimum height=13mm, inner sep=1.5pt},
  io/.style={blk, rounded corners=1.5mm},
  mname/.style={align=left, text width=25.5mm, anchor=east},
  arr/.style={-{Latex[length=1.6mm]}}]
\node[io] (a0) at (0,0) {Ảnh 160$\times$120$\times$3};
\node[blk, right=of a0] (a1) {6 tầng tích chập 3$\times$3 và depthwise, gộp giữa các tầng};
\node[blk, right=of a1] (a2) {TinyFPN gộp 3 mức: stride 8, 16, 32};
\node[blk, right=of a2] (a3) {Mỗi mức: điểm, hộp, 5 điểm mốc};
\node[io, right=of a3] (a4) {Giải mã, NMS bằng C++};
\node[mname] at ([xshift=-2mm]a0.west) {\textbf{YuNet}\\0,07~M tham số\\25,2~M MAC};
\node[io] (b0) at (0,-2.0) {Vùng 80$\times$80$\times$3, cắt 2,7 lần mặt};
\node[blk, right=of b0] (b1) {Tích chập 3$\times$3 bước 2, depthwise};
\node[blk, right=of b1] (b2) {3 nhóm khối dư, khối cuối mỗi nhóm có SE; bản đồ 40, 20, 10, 5};
\node[blk, right=of b2] (b3) {1$\times$1, depthwise 5$\times$5 toàn bản đồ, FC};
\node[io, right=of b3] (b4) {Softmax 3 lớp, lấy xác suất lớp thật};
\node[mname] at ([xshift=-2mm]b0.west) {\textbf{MiniFASNetV1SE}\\0,43~M tham số\\40,9~M MAC};
\node[io] (c0) at (0,-4.0) {Mặt đã căn chỉnh 112$\times$112$\times$3};
\node[blk, right=of c0] (c1) {Tích chập 3$\times$3 bước 2, depthwise, 64 kênh};
\node[blk, right=of c1] (c2) {15 khối cổ chai có ECA; bản đồ 56, 28, 14, 7};
\node[blk, right=of c2] (c3) {1$\times$1 lên 512 kênh, GDC 7$\times$7, FC};
\node[io, right=of c3] (c4) {Vector đặc trưng 512 chiều, chuẩn hóa L2};
\node[mname] at ([xshift=-2mm]c0.west) {\textbf{MobileFaceNet-ECA}\\1,18~M tham số\\221,2~M MAC};
\foreach \r in {a,b,c} {
  \draw[arr] (\r0) -- (\r1);
  \draw[arr] (\r1) -- (\r2);
  \draw[arr] (\r2) -- (\r3);
  \draw[arr] (\r3) -- (\r4);
}
\end{tikzpicture}
\caption{Kiến trúc ba mô hình ở mức khối.}
\label{fig:kien_truc_mo_hinh}
\end{figure}

\section{Thiết kế firmware}
\label{section:3.5}

\subsection{Phân tầng}
\label{section:3.5.1}

Firmware gồm 22 thành phần và \texttt{main}, chia tám tầng như Hình~\ref{fig:tang}. Luật
chính là phụ thuộc chỉ đi xuống, không đi ngang, và CI chặn mọi phụ thuộc sai luật
này. Đây là phân tầng nới: một tầng được gọi thẳng các tầng thấp hơn, không nhất thiết
qua tầng liền kề, và \texttt{common} là cột dùng chung cho mọi tầng. \texttt{ui\_kiosk} và
\texttt{svc\_attendance} không gọi nhau mà gặp qua hàng đợi do \texttt{main} nối. Driver và
tầng mạng viết bằng C. Các thành phần \texttt{svc\_*}, \texttt{ai\_engine} và
\texttt{ui\_kiosk} viết bằng C++, tắt ngoại lệ và RTTI, không cấp phát sau khi khởi động
xong, và header công khai chỉ chứa cú pháp C.

\begin{figure}[!htb]
\centering
\begin{tikzpicture}[
  font=\footnotesize,
  row/.style={draw, align=center, text width=104mm, minimum height=7.6mm, inner sep=3pt},
  tag/.style={anchor=east}]
\node[row] (l7) at (0,0) {\texttt{main}: khởi động, nối dây, tạo task};
\node[row, below=0.8mm of l7] (l6) {\texttt{ui\_kiosk}};
\node[row, below=0.8mm of l6] (l5) {\texttt{svc\_attendance}, \texttt{svc\_sync}};
\node[row, below=0.8mm of l5] (l4) {\texttt{svc\_vision}, \texttt{svc\_door}};
\node[row, below=0.8mm of l4] (l3) {\texttt{ai\_engine}, \texttt{svc\_facedb}, \texttt{drv\_touch}, \texttt{drv\_tof}, \texttt{drv\_audio},\\\texttt{net\_wifi}, \texttt{net\_mqtt}, \texttt{net\_ota}, \texttt{net\_provision}};
\node[row, below=0.8mm of l3] (l2) {\texttt{drv\_camera}, \texttt{drv\_lcd}, \texttt{drv\_servo}, \texttt{drv\_ioexp},\\\texttt{sys\_storage}, \texttt{sys\_time}};
\node[row, below=0.8mm of l2] (l1) {\texttt{bsp\_board}: chân, bus, \texttt{app\_config.h}};
\node[tag] at (l7.west) {L7};
\node[tag] at (l6.west) {L6};
\node[tag] at (l5.west) {L5};
\node[tag] at (l4.west) {L4};
\node[tag] at (l3.west) {L3};
\node[tag] at (l2.west) {L2};
\node[tag] at (l1.west) {L1};
\draw ([xshift=1.5mm]l7.north east) rectangle ([xshift=11mm]l1.south east);
\path ([xshift=6.25mm]l7.north east) -- node[rotate=90] {L0: \texttt{common}} ([xshift=6.25mm]l1.south east);
\end{tikzpicture}
\caption{Các tầng của firmware.}
\label{fig:tang}
\end{figure}

\subsection{Task và đồng bộ}
\label{section:3.5.2}

Bảng~\ref{tab:task} liệt kê các task của ứng dụng. Lõi 1 chỉ chạy task AI: một lượt mô hình
chiếm CPU liên tục, riêng nhận diện 246~ms, nên để chung thì giao diện đứng cả phần tư giây
và Wi-Fi mất gói. Mọi ưu tiên đều dưới 18 để không chèn ngăn xếp Wi-Fi. Task có hạn cứng
như camera, cảm biến, âm thanh đặt cao hơn giao diện và việc nền.

\begin{table}[!htb]
\centering
\small
\caption{Các task của ứng dụng.}
\label{tab:task}
\begin{tabular}{|L{1.8cm}|c|c|L{3.2cm}|L{5.3cm}|}
\hline
\textbf{Task} & \textbf{Lõi} & \textbf{Ưu tiên} & \textbf{Kích hoạt} & \textbf{Việc} \\
\hline
\texttt{cam} & 0 & 7 & Mỗi khung & Lấy khung, vẽ xem trước kèm lớp phủ, đưa khung cho AI \\
\hline
\texttt{tof} & 0 & 6 & Ngắt, hoặc 100~ms & Đọc khoảng cách, báo có người \\
\hline
\texttt{audio} & 0 & 6 & Chờ hàng đợi & Phát tiếng bíp \\
\hline
\texttt{touch} & 0 & 5 & Ngắt GT911 & Đọc điểm chạm \\
\hline
\texttt{ai} & 1 & 5 & Chờ khung & Chạy pipeline \\
\hline
\texttt{ui} & 0 & 4 & 20~ms & Dựng lớp phủ, xử lý chạm \\
\hline
\texttt{attend} & 0 & 4 & Chờ kết quả & Máy trạng thái, ghi lượt, mở thanh chắn \\
\hline
\texttt{mqtt} & 0 & 3 & Thư viện esp-mqtt & Gửi nhận qua TLS \\
\hline
\texttt{ota} & 0 & 3 & Lệnh, mốc đổi JWT & Tải bản cập nhật, xin và đổi JWT \\
\hline
\texttt{net} & 0 & 3 & Một lần lúc khởi động & Chờ Wi-Fi, nối broker, mở SNTP \\
\hline
\texttt{sync} & 0 & 2 & 5~s hoặc có dữ liệu & Gửi bù, nhịp sống, áp lệnh \\
\hline
\end{tabular}
\end{table}

Khung camera đi tới AI qua hàng đợi độ sâu 1, và người gửi lấy khung cũ ra rồi mới gửi, nên
AI luôn xử lý khung mới nhất và trễ không tích lũy. Khung không đi tới task giao diện:
\texttt{cam} vẽ ngay tại chỗ nó đang cầm khung, còn \texttt{ui} chỉ đưa xuống một ảnh lớp
phủ. Chỉ \texttt{cam} ghi vào panel, kể cả các màn phủ kín: panel không cần khóa và không có lượt bàn giao nào gây xé hình. Điểm chạm và phán quyết hiển thị là mức, người đọc chỉ cần giá trị mới nhất; hai thứ này đi bằng biến nguyên tử và ô chung có khóa ngắn.

Hình~\ref{fig:task} vẽ các task nối với nhau ra sao. Mỗi mũi tên liền là một hàng đợi, số trong ngoặc là mức ưu tiên. Task AI đứng một
mình ở lõi 1 và chỉ nhận khung mới nhất. Kết quả của nó, tín hiệu có người từ ToF và lệnh từ máy chủ đều đổ về các task ở lõi 0 qua hàng
đợi. Hai mũi tên đứt là điểm chạm và ảnh lớp phủ, đi qua ô chung vì người đọc chỉ cần giá trị mới nhất.

\begin{figure}[!htb]
\centering
\begin{tikzpicture}[
  font=\footnotesize,
  tk/.style={draw, align=center, text width=19mm, minimum height=8mm, inner sep=1pt},
  q/.style={-{Latex[length=2mm]}},
  sh/.style={densely dashed, -{Latex[length=2mm]}},
  lab/.style={font=\scriptsize, fill=white, inner sep=1pt}]
\draw[densely dashed] (-0.1,2.3) rectangle (9.75,-6.75);
\node[anchor=north west, font=\scriptsize] at (-0.05,2.27) {Lõi 0};
\draw[densely dashed] (10.2,2.3) rectangle (13.2,0.35);
\node[anchor=north west, font=\scriptsize] at (10.25,2.27) {Lõi 1};
\node[tk] (cam) at (1.1,1.3) {\texttt{cam} (7)};
\node[tk] (ai) at (11.7,1.3) {\texttt{ai} (5)};
\node[tk] (tof) at (1.1,-0.4) {\texttt{tof} (6)};
\node[tk] (att) at (8.1,-0.4) {\texttt{attend} (4)};
\node[tk] (touch) at (1.1,-2.1) {\texttt{touch} (5)};
\node[tk] (ui) at (4.6,-2.1) {\texttt{ui} (4)};
\node[tk] (aud) at (8.1,-2.1) {\texttt{audio} (6)};
\node[tk] (net) at (1.1,-4.4) {\texttt{net} (3)};
\node[tk] (mq) at (4.6,-4.4) {\texttt{mqtt} (3)};
\node[tk] (sy) at (8.1,-4.4) {\texttt{sync} (2)};
\node[tk] (ota) at (4.6,-6.0) {\texttt{ota} (3)};
\draw[q] (cam) -- node[lab, above] {\texttt{q\_frame\_ai}, sâu 1} (ai);
\draw[q] (ai.south) |- node[lab, pos=0.75, above] {\texttt{q\_result}} ([yshift=1.5mm]att.east);
\draw[q] (tof) -- node[lab, above] {\texttt{q\_presence}} (att);
\draw[q] (att) -- node[lab, right] {\texttt{q\_audio}} (aud);
\draw[q] (ui) -- node[lab, above] {\texttt{q\_audio}} (aud);
\draw[sh] (touch) -- node[lab, above] {điểm chạm} (ui);
\draw[sh] (ui.north) -- node[lab, sloped, below, pos=0.55] {ảnh lớp phủ} (cam.south east);
\draw[q] ([yshift=-1.5mm]att.east) -- ++(0.4,0) |- node[lab, pos=0.25, right] {\texttt{q\_uplink}} (sy.east);
\draw[q] (ui.south) -- ++(0,-0.55) -| node[lab, pos=0.25, above] {\texttt{q\_roster}} (sy.north);
\draw[q] ([yshift=1.5mm]mq.east) -- node[lab, above, align=center] {\texttt{q\_cmd},\\\texttt{q\_roster}} ([yshift=1.5mm]sy.west);
\draw[q] ([yshift=-1.5mm]sy.west) -- node[lab, below] {gửi qua TLS} ([yshift=-1.5mm]mq.east);
\draw[q] (net) -- node[lab, above] {nối broker} (mq);
\draw[q] (mq) -- node[lab, right] {\texttt{q\_ota}} (ota);
\end{tikzpicture}
\caption{Các task của firmware và hàng đợi nối chúng.}
\label{fig:task}
\end{figure}

Bốn khóa \texttt{m\_facedb\_io}, \texttt{m\_facedb}, \texttt{m\_littlefs}, \texttt{m\_i2c}
luôn lấy theo đúng thứ tự đó, và mọi lần chờ khóa đều có thời hạn. Ghi 576~KB bảng mẫu xuống
flash mất 1,8--2,3~s, lâu hơn hạn 200~ms của lượt tra bảng. Giữ \texttt{m\_facedb} suốt phép
ghi thì lượt tra hết hạn và mặt quen bị báo không nhận ra, nên phép ghi chạy dưới khóa riêng
\texttt{m\_facedb\_io}. Ngắt chỉ gọi các hàm \texttt{FromISR} rồi nhường CPU.

\subsection{Máy trạng thái chấm công}
\label{section:3.5.3}

Nghiệp vụ ở máy là một bảng chuyển trạng thái nằm trong flash, gồm năm trạng thái ở
Hình~\ref{fig:may_trang_thai}. Đồ án cố ý không dùng mẫu State: mẫu đó tốn năm lớp, năm bảng
hàm ảo, và không cho nhìn toàn bộ sơ đồ ở một chỗ. Sự kiện không có trong bảng thì bị bỏ. Các
đường từ \texttt{Idle} và \texttt{Cooldown} thẳng tới \texttt{Granted}, \texttt{Denied} cũng
có mặt, để một phán quyết đến đúng lúc máy đang rảnh không bị nuốt; hình chỉ vẽ các đường
chính.

\begin{figure}[!htb]
\centering
\begin{tikzpicture}[
  font=\footnotesize,
  st/.style={draw, rounded corners, align=center, text width=22mm, minimum height=9mm, inner sep=1.5pt},
  lab/.style={font=\scriptsize, align=center},
  arr/.style={-{Latex[length=2mm]}}]
\node[st] (i) at (0,0) {\texttt{Idle}};
\node[st] (d) at (5.0,0) {\texttt{Detecting}};
\node[st] (g) at (10.4,1.35) {\texttt{Granted}};
\node[st] (n) at (10.4,-1.35) {\texttt{Denied}};
\node[st] (c) at (5.0,-2.7) {\texttt{Cooldown}};
\fill (-1.95,0) circle (1.3mm);
\draw[arr] (-1.82,0) -- (i.west);
\draw[arr] ([yshift=2mm]i.east) -- node[lab, above] {có người, mặt còn xa} ([yshift=2mm]d.west);
\draw[arr] ([yshift=-2mm]d.west) -- node[lab, below] {hết mặt, hết người} ([yshift=-2mm]i.east);
\draw[arr] (d.east) -- node[lab, above, sloped] {khớp} (g.west);
\draw[arr] (d.east) -- node[lab, below, sloped] {giả, chưa nhận ra} (n.west);
\draw[arr] ([xshift=-3mm]g.south) -- node[lab, left] {giả, chưa nhận ra} ([xshift=-3mm]n.north);
\draw[arr] ([xshift=3mm]n.north) -- node[lab, right] {khớp} ([xshift=3mm]g.south);
\draw[arr] (g.east) -- ++(0.55,0) |- node[lab, pos=0.75, below] {hết 2,5~s} ([yshift=-2.5mm]c.east);
\draw[arr] (n.south) |- node[lab, pos=0.75, above] {hết 2~s} ([yshift=2.5mm]c.east);
\draw[arr] (c.west) -| node[lab, pos=0.25, below] {hết 1,5~s, hết người} (i.south);
\end{tikzpicture}
\caption{Máy trạng thái chấm công, các đường chính.}
\label{fig:may_trang_thai}
\end{figure}

Máy cấp quyền theo lần đến của một người. Thiếu luật này thì thanh chắn mở lại và loa kêu lại khoảng mỗi 4~s suốt lúc người đó còn đứng trước máy. Cùng người
trong vòng 5 phút mà chưa rời đi thì máy im lặng: không mở thanh chắn, không ghi, không kêu.
Đã rời khung ngắm rồi quay lại trong 5 phút thì thanh chắn vẫn mở nhưng không sinh bản ghi
mới, vì người ta quét lại thường do thanh chắn chưa kịp mở. Người khác bước tới được cấp ngay
mà không phải chờ thanh chắn của người trước đóng.

\subsection{Bộ nhớ và lưu trữ}
\label{section:3.5.4}

Bảng~\ref{tab:phan_vung} là bảng phân vùng 16~MB. Khe ứng dụng rộng 2,75~MB vì riêng thư viện ESP-DL
thêm 858~KB, đưa ảnh firmware lên 2.679.488~B, không vừa khe 2~MB. Hai khe mô hình tách khỏi
khe ứng dụng nên mô hình cập nhật được độc lập và có đường lùi riêng (F9). Phân vùng \texttt{storage} giữ
một địa chỉ cố định qua mọi bảng phân vùng, nên đổi bảng không làm mất bảng mẫu và nhật ký lượt chấm. Bảng phân vùng
không đi qua OTA được: đổi nó thì phải nạp lại qua USB.

\begin{table}[!htb]
\centering
\small
\caption{Bảng phân vùng flash 16~MB.}
\label{tab:phan_vung}
\begin{tabular}{|L{3.1cm}|c|L{7.9cm}|}
\hline
\textbf{Phân vùng} & \textbf{Cỡ} & \textbf{Chứa} \\
\hline
\texttt{nvs} & 24~KB & Wi-Fi, JWT, cấu hình \\
\hline
\texttt{ota\_0}, \texttt{ota\_1} & 2 $\times$ 2,75~MB & Hai khe firmware \\
\hline
\texttt{models\_0}, \texttt{models\_1} & 2 $\times$ 2,25~MB & Hai khe gói mô hình, đọc qua mmap \\
\hline
\texttt{assets} & 1,5~MB & Âm thanh, chỉ đọc \\
\hline
\texttt{storage} & 4~MB & LittleFS: bảng mẫu, nhật ký lượt chấm \\
\hline
\texttt{coredump} & 64~KB & Bản chụp khi firmware sập \\
\hline
\end{tabular}
\end{table}

Nhật ký lượt chấm chỉ ghi thêm, mỗi bản ghi 48~B có CRC32 (Phụ lục~\ref{appendix:ban_ghi}).
Phân vùng 4~MB chứa khoảng 87.000 bản ghi, thừa cho vài tháng mất mạng. Nhật ký chia thành các
tệp 256~KB, và tệp cũ chỉ bị xóa khi con trỏ gửi đã vượt qua nó. Bản ghi dở dang do mất điện bị
cắt ở lần ghi sau, về đúng mốc 32 + $k \times$ 48 byte, vì để lại vài byte lẻ thì mọi bản ghi
sau đó đọc sai.

Mã lượt \texttt{local\_id} ghép số lần khởi động với số thứ tự trong phiên. Một mã lặp sẽ bị máy
chủ bỏ như bản trùng mà không báo gì, nên số lần khởi động giữ ở NVS và tệp \texttt{boot.bin}. Lúc
khởi động, máy lấy số lớn nhất trong hai chỗ đó và bản ghi cuối của nhật ký rồi cộng một, và máy
chủ còn trả lại mốc lớn nhất khi cấp JWT. Khe còn lại là máy mất NVS đúng lần đầu chạy firmware có
\texttt{boot.bin} trong khi nhật ký đang rỗng: các lượt làm ra trước lần duyệt lại có thể trùng số
đã gửi.

Bảng mẫu giữ tối đa 1.000 bản ghi 576~B, tức khoảng 333 người với ba mẫu mỗi người, nạp lên PSRAM
để quét. Phép ghi đi qua tệp tạm rồi đổi tên, nên mất điện giữa chừng chỉ còn bản cũ hoặc bản mới,
và CRC sai thì rơi về tệp \texttt{.bak}.

\subsection{Đường lượt chấm khi mất mạng}
\label{section:3.5.5}

Hình~\ref{fig:luot_cham} là đường đi của một lượt chấm, thiết kế cho NF5. Bản ghi xuống flash trước
khi có ai biết tới nó, hàng đợi trong RAM chỉ là lời nhắc. Con trỏ đã gửi chỉ tiến sau khi broker
trả PUBACK, nên mất điện giữa chừng làm gửi lại chứ không làm mất, và máy chủ bỏ bản trùng bằng ràng
buộc unique. Đây là giao ít nhất một lần, có chủ ý. API nghe broker bằng một phiên bền và chỉ trả ack sau khi ghi xong. Lượt tới lúc API đang khởi động lại vì thế nằm chờ ở broker. Broker không lưu phiên xuống đĩa: chính EMQX khởi động lại thì mất những tin đang chờ, tối đa 10.000 tin mỗi phiên.

\begin{figure}[!htb]
\centering
\begin{tikzpicture}[font=\footnotesize]
\foreach \x/\n in {0/\texttt{attend}, 2.9/LittleFS, 5.8/\texttt{sync}, 8.7/EMQX, 12.0/Backend} {
  \sqobj{\x}{19mm}{\n}{-6.85}
}
\sqexec{2.9}{-0.9}{-1.25}
\draw[sqsync] (0,-0.9) -- node[sqlab] {ghi 48~B, đồng bộ} (2.82,-0.9);
\draw[sqasync] (0,-1.7) -- node[sqlab, pos=0.75] {nhắc qua \texttt{q\_uplink}} (5.8,-1.7);
\sqexec{2.9}{-2.3}{-2.8}
\draw[sqsync] (5.8,-2.3) -- node[sqlab] {đọc từ con trỏ} (2.98,-2.3);
\draw[sqreply] (2.98,-2.8) -- node[sqlab] {bản ghi} (5.8,-2.8);
\draw[sqasync] (5.8,-3.4) -- node[sqlab] {PUBLISH, QoS 1} (8.7,-3.4);
\draw[sqasync] (8.7,-3.95) -- node[sqlab] {PUBACK} (5.8,-3.95);
\sqexec{2.9}{-4.5}{-4.85}
\draw[sqsync] (5.8,-4.5) -- node[sqlab] {đẩy con trỏ} (2.98,-4.5);
\draw[sqasync] (8.7,-5.1) -- node[sqlab] {chuyển bản tin} (12.0,-5.1);
\sqselfl{12.0}{-5.5}{ghi dòng,\\bỏ bản trùng}
\draw[sqasync] (12.0,-6.5) -- node[sqlab] {PUBACK} (8.7,-6.5);
\end{tikzpicture}
\caption{Đường đi của một lượt chấm từ máy tới máy chủ.}
\label{fig:luot_cham}
\end{figure}

Giờ của lượt chấm lấy từ DS3231 ngay lúc khởi động, rồi được header
\texttt{Date} của API và sau cùng NTP đặt lại, theo thứ tự tin cậy tăng dần. Giờ không đến từ một đồng hồ đã từng được NTP đặt thì bản ghi bật cờ
giờ chưa đồng bộ. Máy chủ không tin một mốc giờ vô lý: lượt mang giờ trước năm 2020 hoặc đi trước giờ
nhận quá một ngày được lưu với cờ nghi vấn và chờ nhân sự sửa. Lượt cũ hơn 7 ngày hay rơi vào kỳ
lương đã chốt thì chờ duyệt mới được tính.

\subsection{Nghỉ, khởi động và cập nhật}
\label{section:3.5.6}

Máy có đúng hai mức (F4). Sau 60~s không có nguồn đánh thức nào, mô hình dừng, đèn nền tắt, màn hình
và camera vào chế độ chờ. VL53L1X thấy người trong 700~mm hay một lần chạm đưa máy về mức thức. Đồ án
bỏ mức giữa vì cho mô hình nghỉ sớm mà màn hình còn sáng là một vùng chết: người dùng thấy máy đang
sống nhưng không mô hình nào nhìn khung hình, trong khi camera 120~mA và đèn nền 100~mA vẫn bật. Mức
điện cắt được chưa đo trên board.

Lúc khởi động, ngoại vi chia hai nhóm. Lưu trữ, màn hình, camera, servo mà hỏng thì dừng, vì không có
chúng thì không chấm công được và cũng không báo được cho người dùng biết vì sao. Mô hình, cảm ứng,
cảm biến, loa, Wi-Fi hỏng thì máy ghi lỗi và chạy tiếp ở mức suy giảm, và tầng trên phải hỏi trạng
thái trước khi dùng.

Bản cập nhật đi bằng lệnh \texttt{down/ota} mang địa chỉ, cỡ và sha256. Địa chỉ là link ký riêng cho
máy đó, hết hạn sau 24 giờ, và máy chỉ tải từ đúng gốc API đã biên dịch sẵn, vì sha256 đi cùng chính
lời mời nên tự nó không chặn được một lời mời giả. Task OTA tải qua HTTPS, ghi vào khe không chạy,
kiểm sha256 rồi đổi khe. Firmware mới phải tự xác nhận hoạt động ở lần khởi động đầu, nếu không
bootloader quay về bản cũ (F9). Bản firmware do CI đưa lên khi số phiên bản tăng, bản mô hình do máy
huấn luyện đưa lên, và người vận hành chỉ bấm chọn máy. Mô hình nhận diện mới đổi không gian vector,
nên sau khi cập nhật mọi người phải đăng ký lại, và bước phát hành đòi xác nhận riêng cho trường hợp
này.

Hình~\ref{fig:td_ota} là trình tự một lần cập nhật firmware. Máy chỉ xác nhận khe mới khi đã chạy đủ 30~s và mô hình đã nạp. Trước đó,
mọi lần khởi động lại đều đưa máy về bản cũ, và máy báo việc quay về đúng một lần ở lần khởi động sau.

\begin{figure}[!htb]
\centering
\begin{tikzpicture}[font=\footnotesize]
\sqactor{0}{Quản trị}{-11.0}
\foreach \x/\w/\n in {3.7/20mm/API, 7.4/20mm/EMQX, 11.1/27mm/Máy chấm công} {
  \sqobj{\x}{\w}{\n}{-11.0}
}
\sqexec{3.7}{-1.0}{-1.9}
\draw[sqsync] (0,-1.0) -- node[sqlab] {chọn bản và máy} (3.62,-1.0);
\draw[sqasync] (3.78,-1.6) -- node[sqlab] {\texttt{down/ota}} (7.4,-1.6);
\draw[sqasync] (7.4,-2.3) -- node[sqlab] {link ký 24 giờ, cỡ, sha256} (11.1,-2.3);
\sqselfl{11.1}{-2.75}{kiểm gốc API,\\cỡ vừa khe}
\sqexec{3.7}{-3.75}{-4.3}
\draw[sqsync] (11.1,-3.75) -- node[sqlab, pos=0.75] {GET qua HTTPS} (3.78,-3.75);
\draw[sqreply] (3.78,-4.3) -- node[sqlab, pos=0.75] {luồng byte} (11.1,-4.3);
\sqselfl{11.1}{-4.75}{ghi khe không chạy,\\băm sha256 khi tải}
\sqselfl{11.1}{-5.75}{khớp thì đổi khe,\\khởi động lại}
\draw[sqasync] (11.1,-6.75) -- node[sqlab, pos=0.7] {nhịp sống qua EMQX,\\đang chạy thử} (3.7,-6.75);
\sqfrag{-0.9}{-7.1}{11.6}{-10.75}{alt}
\node[sqnote, anchor=west] at (-0.15,-7.35) {[đủ 30~s và mô hình đã nạp]};
\sqselfl{11.1}{-7.75}{xác nhận khe}
\draw[sqasync] (11.1,-8.5) -- node[sqlab, pos=0.7] {nhịp sống qua EMQX,\\đã lên} (3.7,-8.5);
\draw[densely dashed] (-0.9,-8.9) -- (11.6,-8.9);
\node[sqnote, anchor=west] at (-0.8,-9.15) {[không xác nhận được]};
\sqselfl{11.1}{-9.5}{bootloader quay\\về bản cũ}
\draw[sqasync] (11.1,-10.45) -- node[sqlab, pos=0.7] {\texttt{up/event}:\\\texttt{OTA\_ROLLED\_BACK}} (3.7,-10.45);
\end{tikzpicture}
\caption{Trình tự cập nhật firmware qua OTA.}
\label{fig:td_ota}
\end{figure}

\section{Thiết kế bản tin MQTT, cơ sở dữ liệu và API}
\label{section:3.6}

\subsection{Bản tin MQTT}
\label{section:3.6.1}

Tám topic ở Bảng~\ref{tab:topic} đều bắt đầu bằng \texttt{kiosk/\{deviceId\}/}. Mọi thứ máy gửi nằm
dưới \texttt{up/}, mọi thứ máy nghe nằm dưới \texttt{down/}, kể cả Last Will, nên luật truy cập chỉ
cần cho máy ghi nhánh \texttt{up/} và đọc nhánh \texttt{down/} của chính nó. Một JWT bị đọc trộm vì
thế không ghi được lệnh giả lên \texttt{down/}. Tài khoản dịch vụ của backend làm ngược lại và bị
cấm ghi lên \texttt{up/}, vì bản tin trên đó là lời khai của máy chấm công và trở thành bằng chứng ai
có mặt lúc mấy giờ.

\begin{table}[!htb]
\centering
\small
\caption{Bảng topic MQTT.}
\label{tab:topic}
\begin{tabular}{|L{3.4cm}|c|c|L{7.0cm}|}
\hline
\textbf{Topic} & \textbf{QoS} & \textbf{Retained} & \textbf{Nội dung} \\
\hline
\texttt{up/attendance} & 1 & Không & Một lượt chấm \\
\hline
\texttt{up/heartbeat} & 0 & Có & Nhịp sống mỗi 30~s: firmware, Wi-Fi, bộ nhớ, số lượt chờ gửi \\
\hline
\texttt{up/event} & 1 & Không & Mặt giả, lỗi phần cứng, mở thanh chắn tay, kết quả lệnh \\
\hline
\texttt{up/status} & 1 & Có & \texttt{online}, hoặc \texttt{offline} do broker phát thay \\
\hline
\texttt{up/enroll} & 1 & Không & Mẫu vừa chụp tại máy \\
\hline
\texttt{down/cmd} & 1 & Không & \texttt{OPEN\_DOOR}, \texttt{REBOOT}, \texttt{SET\_CONFIG} \\
\hline
\texttt{down/enroll} & 1 & Không & Thêm hoặc xóa mẫu \\
\hline
\texttt{down/ota} & 1 & Không & Đề nghị cập nhật \\
\hline
\end{tabular}
\end{table}

Bản tin là JSON theo lược đồ trong \texttt{contracts/schema/}. Bảng~\ref{tab:ban_tin} là các trường
của bản tin lượt chấm. Không có trường nào chứa ảnh (NF4). \texttt{localId} là chuỗi vì số 64 bit
vượt miền số nguyên chính xác của JSON ở phía người đọc.

\begin{table}[!htb]
\centering
\small
\caption{Các trường của bản tin lượt chấm.}
\label{tab:ban_tin}
\begin{tabular}{|L{3.8cm}|L{2.6cm}|L{7.1cm}|}
\hline
\textbf{Trường} & \textbf{Kiểu} & \textbf{Ý nghĩa} \\
\hline
\texttt{deviceId} & Chuỗi & Danh tính máy, trùng với topic \\
\hline
\texttt{localId} & Chuỗi số & Mã lượt, cùng \texttt{deviceId} tạo khóa chống trùng \\
\hline
\texttt{employeeId} & Số nguyên & Mã nhân viên, trong dải \texttt{Int} của máy chủ \\
\hline
\texttt{ts} & Số nguyên & Giờ chấm, mili giây từ epoch, UTC \\
\hline
\texttt{direction} & \texttt{IN}, \texttt{OUT} & Máy luôn gửi \texttt{IN}; vào hay ra do máy chủ suy từ lượt đầu và lượt cuối trong ngày \\
\hline
\texttt{matchScore} & Số thực & Độ tương đồng cosine với mẫu khớp \\
\hline
\texttt{livenessScore} & Số thực & Điểm sống, từ 0 tới 1 \\
\hline
\texttt{modelVersion} & Số nguyên & 16 bit thấp của crc32 tiêu đề gói mô hình đang chạy, khớp bốn chữ số cuối của mã gói trong nhịp sống \\
\hline
\texttt{doorOpened} & Logic & Thanh chắn đã mở \\
\hline
\texttt{capturedOffline} & Logic & Chấm lúc không tới được broker \\
\hline
\texttt{clockUnsynced} & Logic & Giờ chưa đồng bộ NTP, chỉ là gần đúng \\
\hline
\end{tabular}
\end{table}

Danh sách người trên máy hội tụ theo số phiên bản: máy mang số phiên bản danh sách trong nhịp sống,
máy chủ thấy số cũ thì đẩy phần còn thiếu, kể cả các lượt xóa. Áp từng thay đổi thì một lệnh xóa gửi
lúc máy mất mạng sẽ không bao giờ tới, và khuôn mặt người đã nghỉ việc vẫn mở được thanh chắn.

Hình~\ref{fig:td_dang_ky} là trình tự đăng ký một khuôn mặt. Máy chỉ xóa cờ chưa báo của một mẫu khi broker trả PUBACK, nên mẫu chụp lúc
mất mạng được gửi lại khi nối lại. Máy chủ giữ bản mã hóa và đẩy mẫu xuống các máy khác của người đó chạy cùng phiên bản mô hình; người đó không phải đứng chụp lại ở từng máy.

\begin{figure}[!htb]
\centering
\begin{tikzpicture}[font=\footnotesize]
\sqactor{0}{Nhân sự}{-11.45}
\foreach \x/\w/\n in {3.0/19mm/API, 6.0/19mm/EMQX, 9.0/27mm/Máy chấm công} {
  \sqobj{\x}{\w}{\n}{-11.45}
}
\sqactor{12.0}{Người tại máy}{-11.45}
\sqexec{3.0}{-1.0}{-2.85}
\draw[sqsync] (0,-1.0) -- node[sqlab] {gán người vào máy} (2.92,-1.0);
\sqselfr[0]{3.0}{-1.35}{kiểm đồng ý\\dữ liệu mặt}
\draw[sqasync] (3.08,-2.3) -- node[sqlab] {\texttt{down/enroll}} (6.0,-2.3);
\draw[sqreply] (2.92,-2.85) -- node[sqlab] {đã gán} (0,-2.85);
\draw[sqasync] (6.0,-3.3) -- node[sqlab] {thêm vào\\danh sách chờ} (9.0,-3.3);
\draw[sqsync] (12.0,-3.9) -- node[sqlab] {chọn tên} (9.0,-3.9);
\sqfrag{7.75}{-4.25}{12.9}{-6.7}{loop}
\node[sqnote, anchor=west] at (8.6,-4.45) {[3 tư thế: thẳng, trái, phải]};
\draw[sqsync] (12.0,-5.0) -- node[sqlab] {đưa mặt theo tư thế} (9.0,-5.0);
\sqselfr{9.0}{-5.4}{xác thực,\\lấy mẫu}
\draw[sqasync] (9.0,-6.4) -- node[sqlab] {bíp} (12.0,-6.4);
\sqselfr{9.0}{-7.25}{ghi bảng mẫu\\xuống flash}
\draw[sqasync] (9.0,-8.15) -- node[sqlab] {\texttt{up/enroll}, từng mẫu} (6.0,-8.15);
\draw[sqasync] (6.0,-8.7) -- node[sqlab] {PUBACK} (9.0,-8.7);
\draw[sqasync] (6.0,-9.25) -- node[sqlab] {chuyển mẫu} (3.0,-9.25);
\sqselfr{3.0}{-9.65}{mã hóa mẫu,\\đánh dấu đã đăng ký}
\draw[sqasync] (3.0,-11.1) -- node[sqlab] {\texttt{down/enroll} tới\\các máy khác} (6.0,-11.1);
\end{tikzpicture}
\caption{Trình tự đăng ký một khuôn mặt tại máy.}
\label{fig:td_dang_ky}
\end{figure}

Vòng đời máy đi qua \texttt{PENDING}, \texttt{APPROVED}, \texttt{REVOKED}. Thu hồi ghi dòng, xóa cache
của broker rồi đá phiên, theo đúng thứ tự đó. Phải đá phiên vì broker chỉ kiểm JWT lúc nối, nên không đá thì máy đã thu hồi vẫn gửi tin qua phiên cũ. Máy chủ còn chỉ nhận bản tin từ máy \texttt{APPROVED}, và phiên sót lại không ghi được gì.

Hình~\ref{fig:td_duyet_may} là trình tự đưa một máy mới vào hệ thống. Trước khi được duyệt, máy nhận 202 ở mỗi lần hỏi và vẫn chấm công,
giữ lượt trong flash. Broker không tự kiểm JWT mà hỏi API ở mỗi lần nối, nên máy đã bị thu hồi không nối lại được.

\begin{figure}[!htb]
\centering
\begin{tikzpicture}[font=\footnotesize]
\foreach \x/\w/\n in {0/27mm/Máy chấm công, 3.7/20mm/EMQX, 7.4/20mm/API} {
  \sqobj{\x}{\w}{\n}{-9.6}
}
\sqactor{11.1}{Quản trị}{-9.6}
\sqexec{7.4}{-1.25}{-2.45}
\draw[sqsync] (0,-1.25) -- node[sqlab, pos=0.75] {\texttt{POST /devices/register}\\token lô, \texttt{deviceId}, mã 6 số} (7.32,-1.25);
\sqselfr[0]{7.4}{-1.6}{tạo dòng \texttt{PENDING}}
\draw[sqreply] (7.32,-2.45) -- node[sqlab, pos=0.25] {202} (0,-2.45);
\sqselfr{0}{-2.9}{hiện mã nhận máy}
\sqexec{7.4}{-3.85}{-4.25}
\draw[sqsync] (11.1,-3.85) -- node[sqlab] {duyệt: mã đọc trên\\màn hình, tên, vị trí} (7.48,-3.85);
\sqexec{7.4}{-4.75}{-5.35}
\draw[sqsync] (0,-4.75) -- node[sqlab, pos=0.75] {hỏi lại} (7.32,-4.75);
\draw[sqreply] (7.32,-5.35) -- node[sqlab, pos=0.25] {200, JWT 90 ngày} (0,-5.35);
\sqselfr{0}{-5.8}{ghi JWT vào NVS}
\draw[sqasync] (0,-7) -- node[sqlab] {CONNECT qua TLS:\\\texttt{deviceId}, JWT} (3.7,-7);
\sqexec{7.4}{-7.55}{-8.1}
\draw[sqsync] (3.7,-7.55) -- node[sqlab] {\texttt{POST /mqtt/auth}} (7.32,-7.55);
\draw[sqreply] (7.32,-8.1) -- node[sqlab] {cho phép} (3.7,-8.1);
\draw[sqasync] (3.7,-8.65) -- node[sqlab] {CONNACK} (0,-8.65);
\draw[sqasync] (0,-9.2) -- node[sqlab] {\texttt{up/status}: \texttt{online}} (3.7,-9.2);
\end{tikzpicture}
\caption{Trình tự duyệt một máy chấm công mới.}
\label{fig:td_duyet_may}
\end{figure}

\subsection{Cơ sở dữ liệu}
\label{section:3.6.2}

Cơ sở dữ liệu có 60 bảng, 46 kiểu liệt kê và 102 migration. Hình~\ref{fig:er} vẽ các bảng lõi theo ký hiệu chân chim: đầu ba nhánh ở phía nhiều, hai vạch là đúng một,
vòng tròn là có thể không có. Mọi
phân hệ gắn vào \texttt{Employee} và không phân hệ nào giữ bản sao hồ sơ. Bên trái là phần máy chấm
công và lượt chấm, ở giữa là tổ chức và thời gian, bên phải là lương.

\begin{figure}[!htb]
\centering
\begin{tikzpicture}[
  font=\scriptsize,
  e/.style={draw, align=center, text width=27mm, minimum height=6.5mm, inner sep=1pt},
  k/.style={e, thick, fill=black!8},
  y=1.16cm]
\node[e] (dev) at (0,0) {Device};
\node[e] (enr) at (0,-1.25) {DeviceEnrollment};
\node[e] (rec) at (0,-2.5) {AttendanceRecord};
\node[e] (ftp) at (0,-3.75) {FaceTemplate};
\node[e] (day) at (0,-5.0) {AttendanceDay};
\node[e] (leg) at (3.75,0) {LegalEntity};
\node[e] (dep) at (3.75,-1.25) {Department};
\node[k] (emp) at (3.75,-2.5) {Employee};
\node[e] (req) at (3.75,-3.75) {Request};
\node[e] (bal) at (3.75,-5.0) {LeaveBalance};
\node[e] (lty) at (3.75,-6.25) {LeaveType};
\node[e] (usr) at (7.5,0) {User};
\node[e] (cmp) at (7.5,-1.25) {CompensationRecord};
\node[e] (asg) at (7.5,-2.5) {ShiftAssignment};
\node[e] (shf) at (7.5,-3.75) {Shift};
\node[e] (pol) at (11.25,0) {PayrollPolicy};
\node[e] (tax) at (11.25,-1.25) {TaxBracket};
\node[e] (per) at (11.25,-2.5) {PayrollPeriod};
\node[e] (run) at (11.25,-3.75) {PayrollRun};
\node[e] (psl) at (11.25,-5.0) {Payslip};
\node[e] (lin) at (11.25,-6.25) {PayslipLine};
\draw[erfk] (enr) -- (dev);
\draw[erfk] (enr.south east) -- (emp.north west);
\draw[erfk] (rec.north west) to[bend left=35] (dev.south west);
\draw[erfk] (rec) -- (emp);
\draw[erfk] (ftp.north east) -- (emp.south west);
\draw[erfk] (day.north east) -- ([xshift=4mm]emp.south west);
\draw[erfk] (dep) -- (leg);
\draw[erfkopt] (emp) -- (dep);
\draw[eroo] (usr.south west) -- (emp.north east);
\draw[erfk] (cmp.south west) -- ([xshift=-4mm]emp.north east);
\draw[erfk] (asg) -- (emp);
\draw[erfk] (asg) -- (shf);
\draw[erfk] (req) -- (emp);
\draw[erfkopt] (req.west) -- ++(-0.3,0) |- (lty.west);
\draw[erfk] (bal) -- (lty);
\draw[erfk] (bal.east) -- ++(0.2,0) |- ([yshift=-0.7mm]emp.east);
\draw[erfk] (psl.west) -| ([xshift=4mm,yshift=-1.9mm]emp.east) -- ([yshift=-1.9mm]emp.east);
\draw[erfk] (tax) -- (pol);
\draw[erfkopt] (pol.north) -- ++(0,0.3) -| ([xshift=3mm]leg.north);
\draw[erfk] (per.east) -- ++(0.3,0) -- ++(0,3.4) -| ([xshift=-3mm]leg.north);
\draw[erfk] (run) -- (per);
\draw[erfk] (psl) -- (run);
\draw[erfk] (lin) -- (psl);
\end{tikzpicture}
\caption{Các bảng lõi và quan hệ chính, ký hiệu chân chim.}
\label{fig:er}
\end{figure}

Năm quyết định chi phối các bảng này: (i)~dữ liệu có thời hạn không bị sửa đè,
\texttt{CompensationRecord} và \texttt{PayrollPolicy} mang ngày hiệu lực, tăng lương là thêm một dòng và phiếu tháng cũ vẫn tính được đúng; (ii)~lượt chấm thô tách khỏi ngày công,
\texttt{AttendanceRecord} là sự kiện còn \texttt{AttendanceDay} là kết quả, mỗi người mỗi ngày một
dòng; (iii)~phiếu lương chép lại những gì nó dựa vào, gồm tài khoản ngân hàng và dòng chính sách đã
dùng, nên đổi hồ sơ về sau không đổi phiếu đã phát; (iv)~luật mà vi phạm là sai công hay sai tiền thì đặt thành ràng buộc của cơ sở dữ liệu, gồm
ràng buộc unique trên cặp máy và mã lượt, ràng buộc loại trừ chặn đơn nghỉ chồng ngày, và \texttt{CHECK}
cho số ngày phép và số tiền trên dòng phiếu; (v)~tiền lưu bằng số nguyên đồng, tỉ lệ lưu bằng phần vạn.
\texttt{PayslipLine} giữ mỗi khoản một dòng, để câu hỏi như quý này trả bao nhiêu tiền tăng ca trả lời
được bằng một phép gộp SQL.

\subsection{API và phân quyền}
\label{section:3.6.3}

API theo kiểu REST, 260 đường trong 27 module, gom nhóm ở Bảng~\ref{tab:api}. Mọi đường có tài liệu
Swagger và trả lỗi bằng mã, frontend mới đổi mã thành câu theo ngôn ngữ người dùng (NF8). Danh sách
phân trang bằng con trỏ, và số đếm có trần để một phép đếm không quét cả bảng (NF7). Việc nặng như tính
lương, dựng bảng công, gửi thông báo đi qua hàng đợi và phải chạy lại được mà không đổi kết quả.
Trang đang mở nhận tin thay đổi qua Socket.IO, nên người dùng không phải tải lại.

\begin{table}[!htb]
\centering
\small
\caption{Các nhóm API.}
\label{tab:api}
\begin{tabular}{|L{3.8cm}|L{7.9cm}|c|}
\hline
\textbf{Nhóm} & \textbf{Module} & \textbf{Số đường} \\
\hline
Xác thực, tài khoản & \texttt{auth}, \texttt{users} & 23 \\
\hline
Tổ chức và hồ sơ & \texttt{employees}, \texttt{org}, \texttt{profile}, \texttt{certificates}, \texttt{documents}, \texttt{onboarding}, \texttt{assets} & 81 \\
\hline
Thời gian & \texttt{attendance}, \texttt{timesheet}, \texttt{shifts}, \texttt{leave} & 40 \\
\hline
Lương & \texttt{payroll}, \texttt{compensation}, \texttt{policy}, \texttt{disputes} & 52 \\
\hline
Máy chấm công & \texttt{devices}, \texttt{enrollment}, \texttt{models} & 32 \\
\hline
Báo cáo, tiện ích & \texttt{reports}, \texttt{notifications}, \texttt{search}, \texttt{audit}, \texttt{health} & 32 \\
\hline
\end{tabular}
\end{table}

Phân quyền tính tới từng dòng dữ liệu: cùng một đường API trả về tập dòng khác nhau theo người gọi
(NF6). Bảng~\ref{tab:quyen} tóm tắt phạm vi của từng vai. Mọi đường đọc đi qua một dịch vụ phạm vi
chung, và đường có mã đối tượng trả 404 khi đối tượng nằm ngoài phạm vi, để người dò không biết đối
tượng có tồn tại hay không. Tài khoản nhân sự và lương có thể giới hạn theo pháp nhân. Không ai tự
duyệt đơn hay tự sửa lương của mình, và mọi thao tác ghi để lại nhật ký.

\begin{table}[!htb]
\centering
\small
\caption{Phạm vi dữ liệu theo vai.}
\label{tab:quyen}
\begin{tabular}{|L{2.3cm}|L{6.0cm}|L{5.2cm}|}
\hline
\textbf{Vai} & \textbf{Thấy} & \textbf{Sửa} \\
\hline
\texttt{VIEWER} & Chỉ chính mình, vai mặc định & Không gì \\
\hline
\texttt{EMPLOYEE} & Hồ sơ, công, phép, phiếu của mình & Đơn của mình, vài ô liên lạc \\
\hline
\texttt{MANAGER} & Thêm cây dưới quyền, không có lương và giấy tờ riêng & Duyệt đơn của cây \\
\hline
\texttt{HR} & Toàn bộ hồ sơ, công, phép trong pháp nhân phụ trách & Hồ sơ, hợp đồng, phép, phân ca \\
\hline
\texttt{PAYROLL} & Như \texttt{HR}, thêm lương và phiếu lương & Chạy kỳ lương, chốt kỳ \\
\hline
\texttt{ADMIN} & Tất cả, thêm máy chấm công và tài khoản & Tất cả \\
\hline
\end{tabular}
\end{table}

\section{Thiết kế giao diện}
\label{section:3.7}

\subsection{Giao diện máy chấm công}
\label{section:3.7.1}

Máy có chín màn hình nối với nhau như Hình~\ref{fig:man_hinh}. Màn quét là màn chính và luôn có khung
xem trước. Menu có đúng ba dòng Thêm người, Danh sách, Cài đặt, là ba việc người vận hành mở ra để làm,
và Wi-Fi nằm trong Cài đặt. Màn hình không có nút xóa người hay chụp lại, vì menu không khóa và ai đứng
trước máy cũng chạm được; hai việc đó chỉ làm từ dashboard bởi tài khoản đã qua xác thực hai lớp.

\begin{figure}[!htb]
\centering
\begin{tikzpicture}[
  font=\footnotesize,
  node distance=6mm and 8mm,
  s/.style={draw, rounded corners, align=center, text width=24mm, minimum height=9mm, inner sep=1.5pt},
  arr/.style={-{Latex[length=2mm]}}]
\node[s] (scan) {Quét};
\node[s, right=of scan] (menu) {Menu};
\node[s, right=of menu] (enrol) {Thêm người};
\node[s, right=of enrol] (cap) {Chụp mẫu};
\node[s, below=of enrol] (ppl) {Danh sách};
\node[s, below=of ppl] (set) {Cài đặt};
\node[s, right=of set, yshift=7mm] (wifi) {Wi-Fi};
\node[s, right=of set, yshift=-7mm] (devi) {Thiết bị của tôi};
\node[s, below=of scan] (upd) {Cập nhật};
\draw[arr] (scan) -- (menu);
\draw[arr] (menu) -- (enrol);
\draw[arr] (enrol) -- (cap);
\draw[arr] (menu.south) |- (ppl.west);
\draw[arr] (menu.south) |- (set.west);
\draw[arr] (set.east) -- (wifi.west);
\draw[arr] (set.east) -- (devi.west);
\draw[arr] (scan) -- node[font=\scriptsize, right] {khi cập nhật} (upd);
\end{tikzpicture}
\caption{Bản đồ màn hình của máy chấm công.}
\label{fig:man_hinh}
\end{figure}

Khi một màn của người vận hành đang mở, mô hình nghỉ và chấm công dừng hẳn, vì người thao tác đứng ngay
trước camera và để chấm công chạy ngầm là chấm chính họ. Giá phải trả là không ai chấm công được ở máy đó
trong lúc ấy, nên các màn này tự về màn quét sau 30~s không chạm.

Máy tự vẽ giao diện bằng một \texttt{Canvas}, không dùng thư viện LVGL. Màn hình chính là video xem trước,
và thêm LVGL là có hai bộ vẽ cùng ghi một panel SPI. Vì vậy nút, danh sách, thanh trượt, bàn phím đều
tự viết, và không có gì tự chặn chữ tràn khỏi panel 320~px, nên mọi màn bắt buộc vẽ qua tầng \texttt{theme}, giữ bảng màu, thang chữ, nấc giãn cách,
và tầng \texttt{widgets}, giữ các phần tử lặp lại cùng phép xếp dọc. Chuỗi hiển thị khai bằng \texttt{enum class StrId}, và thiếu một bản dịch là lỗi biên dịch. Tiếng Việt là mặc định.

Khung ngắm vẽ theo chính cổng cỡ mặt của pipeline. Đường xem trước phóng khung camera lên panel với hệ số 1,502: cổng 100~px ứng với 150~px trên panel. Người ta lấp khung bằng cả cái đầu, mà tỉ lệ đầu trên mặt là 1,48, nên khung rộng 240 $\times$ 296~px: đầu lấp kín khung cho hộp mặt 108~px trong khung camera,
dư 8~px trên cổng. Khung trắng khi chờ, hổ phách khi cần người dùng sửa, xanh dương khi pipeline đang làm
việc, xanh lá khi đạt. Dòng nhắc dưới khung đi theo đúng các cổng: Đưa khuôn mặt vào khung, Lại gần hơn, Lùi
lại một chút, Đang nhận diện. Kết quả đạt hiện thành thẻ 1,5~s. Máy chỉ phát tiếng khi một việc đã xong là
mở thanh chắn hoặc giữ được một mẫu đăng ký, vì người đang quay đầu không đọc được màn hình.

\subsection{Giao diện web}
\label{section:3.7.2}

Web chia màn hình theo việc người ta đang làm, tránh một thanh bên hai mươi mục ngang hàng theo tên phân
hệ. Việc cần làm nằm ở trang chủ và mang số đếm trên thanh bên, tra cứu nằm ở thanh
bên chính, cấu hình nằm sau mục Thiết lập, còn báo cáo đứng riêng. Một bản Next.js phục vụ mọi vai với 46 trang,
và trang chủ đổi theo vai như Bảng~\ref{tab:trang_chu}.

\begin{table}[!htb]
\centering
\small
\caption{Trang chủ theo vai.}
\label{tab:trang_chu}
\begin{tabular}{|L{3.0cm}|L{2.9cm}|L{7.6cm}|}
\hline
\textbf{Vai} & \textbf{Trang chủ} & \textbf{Nội dung chính} \\
\hline
Nhân viên & Trang của tôi & Hôm nay: ca, các lượt chấm, đúng giờ hay muộn; việc cần làm \\
\hline
Quản lý & Trang của tôi & Thêm việc chờ duyệt và tình hình nhóm hôm nay \\
\hline
Nhân sự, quản trị & Tổng quan & Việc chờ duyệt, hợp đồng và thử việc sắp hết, ngày lệch giờ \\
\hline
Kế toán lương & Tổng quan & Kỳ đang mở theo năm bước: chạy, soát, chốt, trả, gửi phiếu \\
\hline
\end{tabular}
\end{table}

Mỗi trang gồm năm tầng từ trên xuống: thanh bên, thanh trên, thanh tab phụ, đầu trang, thân trang. Nội dung
rộng tối đa 1.400~px. Ở các danh sách, cả dòng mở được bản ghi, bộ lọc nằm ngay trên bảng và giữ trong URL,
số trên trang chủ bấm được và dẫn tới đúng danh sách có số đó. Hành động không đảo ngược được như chốt kỳ
phải xác nhận, còn việc lặp lại nhiều như duyệt đơn thì không.

Thiết kế giả định công nhân, nhân viên kho, lái xe, bảo vệ chỉ có điện thoại, nên cổng nhân viên dựng riêng
cho màn hẹp (F14). Thanh bên đổi thành thanh tab dưới đáy với đúng năm ô: Trang của tôi, Lịch ca, Công của tôi,
Đơn từ, Phiếu lương. Vai có thêm việc đổi ô cuối thành Thêm. Bảng đổi thành thẻ, bộ lọc đổi thành tấm trượt từ đáy, hành động chính neo ở đáy trong tầm ngón
cái. Điện thoại cố ý không chấm công được: vị trí GPS giả được bằng phần mềm, và thêm đường đó là bỏ qua chống
giả mạo và so mẫu khuôn mặt của máy chấm công.

Mạng yếu được coi là trạng thái thường (F15). Dữ liệu đã xem nằm trong ngăn riêng của tài khoản và
bị xóa khi đăng xuất, đơn gửi lúc mất mạng xếp hàng và báo rõ là đang chờ gửi. Trên Android, ứng
dụng là một vỏ Trusted Web Activity~\cite{twa} mở đúng trang web đó, vì Chrome không cho PWA đã cài
tô hai thanh hệ thống theo màu của ứng dụng.

Chương này chốt 15 yêu cầu chức năng và 9 yêu cầu phi chức năng. Phần cứng là mạch đế 158 $\times$ 121~mm với
hai đường 5~V. Firmware có 22 thành phần và \texttt{main} trên tám tầng, 11 task, lõi 1 dành riêng cho AI, và
lượt chấm đi qua flash trước khi lên MQTT. Máy chủ có 60 bảng và 260 đường API.
Chương~\ref{chapter:Implementation} trình bày cách hiện thực các thiết kế này.

\end{document}
